% !TeX root = ../../Book3.tex
% 纲要标题：## 第3章　曲率张量与局部几何的展开
% 纲要位置：几何学计划纲要.md，第 2192 行。
% 正文状态：初稿；尚未正式审核。

% 纲要内容开始
% 
% > 本章目标与证明范围：从协变微分的非交换性得到代数曲率张量，再把各次收缩和法坐标展开联系到可观测的长度、面积和体积变化。
% 
% 纲要内容结束

\chapter{曲率张量与局部几何的展开}
\label{g3:ch:03}
\BookChapterNumber{g3:ch:03}

\begin{chapterabstract}
	法坐标消去度量在中心的一阶变化，二阶变化则由曲率控制。本章从协变微分的非交换性证明曲率的张量性、代数对称性及两条 Bianchi 恒等式，建立曲率算子、截面曲率、Ricci 曲率与标量曲率之间的关系。随后推导联络、度量和体积密度的法坐标展开，把不同收缩解释为长度、面积和小球体积的低阶变化。空间形式、乘积、翘曲积与共形度量提供计算模型，Weyl 分解则分离标量、无迹 Ricci 和无迹四阶曲率三部分。
\end{chapterabstract}

\begin{learninggoals}
	\begin{itemize}
		\item 在向量曲率、四阶曲率张量和二向量曲率算子之间转换，并核对指标次序与球面的正曲率符号。
		\item 证明两条 Bianchi 恒等式及其收缩，说明截面、Ricci、标量和曲率算子正性之间的含义与差别。
		\item 从法坐标径向条件求出联络和度量的低阶项，推导体积密度、小球体积及面积的展开，并控制余项。
		\item 计算乘积、翘曲积及共形变化的曲率，使用 Weyl 分解解释低维确定性和高维信息损失。
	\end{itemize}
\end{learninggoals}

在平面上画一个半径很小的圆，它的周长是 \(2\pi r\)。在半径 \(a\) 的球面上，以沿球面的距离 \(r\) 为半径画圆，周长却是 \(2\pi a\sin(r/a)\)。当 \(r\) 很小时，两者的一阶项相同，差异从三次项开始：
\[
	2\pi a\sin(r/a)=2\pi r-\frac{\pi}{3a^2}r^3+O(r^5).
\]
这个差异可以由长度测量得到，不需要先把球面与环境空间中的某个法向比较。高维中，不同二维方向可以出现不同的差异；将它们沿某个方向相加，会得到 Ricci 曲率，再把方向全部平均，则出现标量曲率。

\ProseParagraphBreak

第二章的法坐标使 \(g_{ij}(p)=\delta_{ij}\)、\(\partial_kg_{ij}(p)=0\)。因此我们可以把度量与 Euclid 度量的首个差异集中在二阶项中。证明将沿两个方向进行：先由联络求导的交换差得到曲率，找出它必须满足的对称性；再由这些曲率分量恢复法坐标中的二次项，计算可观测的几何量。曲率约束之间的区别也将在实际局部度量中检验。

\ProseParagraphBreak

本章使用第一章的 Levi-Civita 联络、体积密度及非负 Laplace 算子 \(\Delta_g=-\operatorname{div}_g\operatorname{grad}_g\)，以及第二章的指数映射、Gauss 引理和法坐标。外积空间只需有限维线性代数：\(X\wedge Y\) 表示由两个向量张成的有向面积元素。流形及度量均光滑，通常假设无边界；紧性、连通性和维数限制在实际需要的结果中另行写明。法坐标展开由坐标恒等式和 Taylor 公式推出。

\ProseParagraphBreak

作为阅读对照，曲率、曲率算子及各次收缩可参看 Petersen 第3版第3.1.1--3.1.6节；球面、乘积及旋转型度量见第4.2.1--4.2.3节，共形表示见第4.3.3节，指数坐标性质见第5.5.3节\cite[第3.1.1--3.1.6节; 第4.2.1--4.2.3节; 第4.3.3节; 第5.5.3节]{Petersen2016}。本章固定的四阶张量次序是 \(\mathrm{Rm}(X,Y,Z,W)=g(R(X,Y)W,Z)\)，因此 \(\mathrm{Rm}(u,v,u,v)\) 是正交单位向量所张截面的曲率。

% !TeX root = ../../Book3.tex
% 纲要标题：### 3.1 曲率的代数结构
% 纲要位置：几何学计划纲要.md，第 2196 行。
% 正文状态：未撰写；本文件目前是框架。

% 纲要内容开始
% 
% * 曲率张量性及符号
% * 四阶张量的对称性
% * 第一、第二 Bianchi 恒等式
% * 曲率算子及二形式上的作用
% 
% 纲要内容结束

\section{曲率的代数结构}
\label{g3:sec:0301}

待撰写

\subsection{曲率张量性及符号}
\label{g3:subsec:0301:01}
待撰写

\subsection{四阶张量的对称性}
\label{g3:subsec:0301:02}
待撰写

\subsection{第一、第二 Bianchi 恒等式}
\label{g3:subsec:0301:03}
待撰写

\subsection{曲率算子及二形式上的作用}
\label{g3:subsec:0301:04}
待撰写
% 纲要条目排为编号小节标题，并进入目录与书签；说明段仍保留于纲要注释中。

% !TeX root = ../../Book3.tex
% 纲要标题：### 3.2 不同曲率约束
% 纲要位置：几何学计划纲要.md，第 2203 行。
% 正文状态：初稿；尚未正式审核。

% 纲要内容开始
% 
% * 截面曲率和二维截面
% * Ricci 曲率与体积效应
% * 标量曲率与收缩
% * 高维不同正性条件的严格区别
% 
% 纲要内容结束

\section{不同曲率约束}
\label{g3:sec:0302}

曲率张量同时接受四个方向，实际几何问题却常常只指定一个二维平面、一个径向方向，或者一小块区域的总体积。截面曲率、Ricci 曲率与标量曲率分别对应这些不同的观察方式。\footnote{Gregorio Ricci-Curbastro（1853--1925），意大利数学家，创立绝对微分学，并与 Levi-Civita 合作发展其几何应用。}它们都来自同一个曲率张量，但取值所保留的信息和“为正”的含义并不相同。本节先建立这些量之间的精确关系，再用可以直接核算的例子辨明各种正性条件。

\subsection{截面曲率和二维截面}
\label{g3:subsec:0302:01}

回顾\cref{g3:def:sectional-curvature}。对 \(\sigma\subset\mathrm{T}_pM\) 的任意一组基 \(X,Y\)，本章的四阶张量约定将截面曲率写为
\begin{equation}
	\label{g3:eq:sectional-bivector}
	K_p(\sigma)
	=\frac{\mathrm{Rm}(X,Y,X,Y)}{|X\wedge Y|_g^2}
	=\frac{\langle\mathcal R(X\wedge Y),X\wedge Y\rangle}{|X\wedge Y|_g^2},
\end{equation}
其中 \(|X\wedge Y|_g^2=g(X,X)g(Y,Y)-g(X,Y)^2\) 是平行四边形面积的平方。因此当 \(u,v\) 为 \(\sigma\) 的正交单位基时，\(K_p(\sigma)=\mathrm{Rm}(u,v,u,v)\)。

\ProseParagraphBreak

这个商与平面中的基无关。事实上，若 \(X'=aX+bY\)、\(Y'=cX+dY\)，则 \(X'\wedge Y'=(ad-bc)X\wedge Y\)。\eqref{g3:eq:sectional-bivector}的分子和分母同时乘以 \((ad-bc)^2\)。特别地，反转平面定向也不改变曲率。这个检验解释了为何必须除以面积平方：仅有四阶张量的分子，会随着所选基的长度和夹角变化。

\begin{definition}
	\label{g3:def:sectional-curvature-algebraic}
	设 \((V,g)\) 为维数 \(n\geq2\) 的实内积空间，\(A\) 为 \(V\) 上的代数曲率张量。对二维子空间 \(\sigma=\operatorname{span}\{X,Y\}\)，定义它的截面曲率为
	\[
		K_A(\sigma)=\frac{A(X,Y,X,Y)}{|X\wedge Y|_g^2}.
	\]
	同样的换基计算保证它只依赖 \(\sigma\)。
\end{definition}

截面曲率看似只读取了四阶张量中两个输入重复的值，却能恢复全部分量。下面的恢复过程除了利用双线性，还将用到第一 Bianchi 恒等式。

\begin{proposition}
	\label{g3:prop:sectional-determines-curvature}
	设 \((V,g)\) 为有限维实内积空间，\(\dim V\geq2\)，\(A,B\) 是其上的代数曲率张量。若它们在每个二维子空间上的截面曲率相同，则 \(A=B\)。
\end{proposition}
\begin{proof}
	令 \(D=A-B\)。截面曲率相等说明对线性无关的 \(x,y\)，有 \(D(x,y,x,y)=0\)；线性相关时，同一式子由反对称性成立。因此这个等式对所有 \(x,y\) 都成立。下面作两次极化，再用第一 Bianchi 恒等式消去剩余分量。

	将第一个重复向量替换为 \(x+z\)，利用交换两对输入的对称性，得到
	\[
		0=D(x+z,y,x+z,y)=2D(x,y,z,y).
	\]
	再在刚得到的等式中将 \(y\) 替换为 \(y+w\)，便有
	\[
		D(x,y,z,w)+D(x,w,z,y)=0.
	\]
	由最后两个输入的反对称性，这等价于
	\begin{equation}
		\label{g3:eq:sectional-zero-polarization}
		D(x,y,z,w)=D(x,w,y,z).
	\end{equation}
	于是交换两对输入再用该式，得到
	\[
		D(y,z,x,w)=D(x,w,y,z)=D(x,y,z,w).
	\]
	同样，
	\[
		\begin{aligned}
			D(z,x,y,w)&=D(y,w,z,x)=D(y,x,w,z)\\
			&=D(x,y,z,w),
		\end{aligned}
	\]
	最后一步同时交换第一对和第二对输入。第一 Bianchi 恒等式的三个循环项因而相等：
	\[
		0=D(x,y,z,w)+D(y,z,x,w)+D(z,x,y,w)
		=3D(x,y,z,w).
	\]
	故 \(D=0\)。
\end{proof}

\begin{corollary}
	\label{g3:cor:pointwise-isotropic-curvature}
	设 \((V,g)\) 为维数至少为二的实内积空间。代数曲率张量 \(A\) 在所有二维子空间上的截面曲率均为同一个实数 \(k\)，当且仅当
	\begin{equation}
		\label{g3:eq:constant-sectional-algebraic}
		A(X,Y,Z,W)
		=k\bigl(g(X,Z)g(Y,W)-g(X,W)g(Y,Z)\bigr).
	\end{equation}
\end{corollary}
\begin{proof}
	右端具有两对输入的反对称性与交换两对输入的对称性，第一 Bianchi 循环和中的六个乘积逐对抵消，故它是代数曲率张量。令 \(Z=X,W=Y\)，分子恰为 \(k|X\wedge Y|_g^2\)，所以它的截面曲率恒为 \(k\)。反过来，若 \(A\) 也有这些截面值，\cref{g3:prop:sectional-determines-curvature}使两个张量相同。
\end{proof}

这里的 \(k\) 只描述一个切空间中的各个平面。若在流形上逐点应用这一推论，首先得到的仍是可能随点变化的 \(k(p)\)；它何时必须为常数，将由本节稍后的微分恒等式决定。

\subsection{Ricci 曲率与体积效应}
\label{g3:subsec:0302:02}

固定一个方向 \(v\)，可以分别观察包含 \(v\) 的各个二维平面。若把正交于 \(v\) 的所有方向同时考虑，对相应曲率求和便比单个截面更适合描述横向体积的变化。这个求和应由张量的迹来定义，使结果不依赖所选的正交基。

\begin{definition}
	\label{g3:def:ricci-scalar}
	设 \((M^n,g)\) 为光滑黎曼流形，\(\mathrm{Rm}\) 采用\cref{g3:def:curvature-four-tensor}的约定。在任意正交单位基 \(e_1,\ldots,e_n\) 下，定义\term[ricciqulv]{Ricci 曲率}[Ricci curvature]为二阶协变张量
	\begin{equation}
		\label{g3:eq:ricci-definition}
		\operatorname{Ric}(Y,W)
		=\sum_{i=1}^n\mathrm{Rm}(e_i,Y,e_i,W)
		=\operatorname{tr}\bigl(X\mapsto R(X,Y)W\bigr).
	\end{equation}
	它的迹称为\term[biaoliangqulv]{标量曲率}[scalar curvature]，记为
	\[
		\operatorname{Scal}=\operatorname{tr}_g\operatorname{Ric}
		=\sum_{j=1}^n\operatorname{Ric}(e_j,e_j).
	\]
\end{definition}

定义中的 \(X\mapsto R(X,Y)W\) 是同一个切空间上的线性自映射，其迹不依赖基。一般坐标中的公式为
\begin{equation}
	\label{g3:eq:ricci-scalar-contractions}
	\operatorname{Ric}_{jl}=g^{ik}\mathrm{Rm}_{ijkl},\qquad
	\operatorname{Scal}=g^{jl}\operatorname{Ric}_{jl}.
\end{equation}
这些公式还说明两者随基点光滑变化。\(\operatorname{Ric}\) 对两个向量线性，\(\operatorname{Scal}\) 则是每一点处的实数。

\symidx{\(\operatorname{Ric}\)}
\symidx{\(\operatorname{Scal}\)}

\begin{proposition}
	\label{g3:prop:ricci-sectional-trace}
	在 \(n\) 维光滑黎曼流形上，\(\operatorname{Ric}\) 是对称二阶张量。若 \(n\geq2\)、\(|v|_g=1\)，并将 \(v=e_1\) 补成正交单位基，则
	\begin{equation}
		\label{g3:eq:ricci-sectional-sum}
		\operatorname{Ric}(v,v)=\sum_{a=2}^nK(\operatorname{span}\{v,e_a\}).
	\end{equation}
	特别地，若一点处所有截面曲率位于 \([a,b]\)，则在该点有二次型不等式
	\[
		(n-1)a\,g\leq\operatorname{Ric}\leq(n-1)b\,g.
	\]
\end{proposition}
\begin{proof}
	由交换两对输入的对称性，
	\[
		\mathrm{Rm}(e_i,Y,e_i,W)=\mathrm{Rm}(e_i,W,e_i,Y),
	\]
	求和即得 Ricci 对称性。将 \(Y=W=v=e_1\) 代入定义，\(i=1\) 的项因反对称性为零，其余每项是一个正交单位二维平面的截面曲率，得到\eqref{g3:eq:ricci-sectional-sum}。最后对单位向量应用逐项上下界，再利用二次齐次性恢复一般向量，即得二次型估计。
\end{proof}

因此 \(\operatorname{Ric}(v,v)/(n-1)\) 是沿单位方向 \(v\) 所见的这些正交截面的平均曲率。求和中的个别截面值可以变化，和却不受补全正交基的方式影响。在维数二，只有一个横向方向，所以 \(\operatorname{Ric}=Kg\)；维数增加后，正负截面可能在同一个方向的求和中抵消。

\ProseParagraphBreak

这个迹如何进入体积，可以先由下一节将证明的展开辨认。固定 \(p\)，在 \(\mathrm{T}_pM\) 中取正交单位基，并用指数映射得到法坐标。对单位向量 \(v\)，\cref{g3:prop:normal-volume-density}将证明
\begin{equation}
	\label{g3:eq:ricci-volume-direction}
	\sqrt{\det(g_{ij}(tv))}
	=1-\frac16\operatorname{Ric}_p(v,v)t^2+O(t^3).
\end{equation}
左侧是指数坐标下的体积密度与切空间 Euclid 体积密度之比。横向各方向的一阶偏差为零，二阶偏差在行列式中相加，因而出现截面曲率的迹。对某个 \(v\) 有 \(\operatorname{Ric}(v,v)>0\) 时，这个方向上足够近的体积比小于一；负值则使二阶修正朝相反方向变化。公式的推导放在法坐标二阶展开之后，当前 Ricci 张量的定义和对称性只用了代数收缩。

\subsection{标量曲率与收缩}
\label{g3:subsec:0302:03}

再对 Ricci 张量取迹，就把对某个方向的观察合成一个数。对任意正交单位基，\eqref{g3:eq:ricci-sectional-sum}给出
\begin{equation}
	\label{g3:eq:scalar-sectional-average}
	\operatorname{Scal}
	=2\sum_{i<j}K(\operatorname{span}\{e_i,e_j\})
	=2\operatorname{tr}\mathcal R.
\end{equation}
每个坐标二维平面在 Ricci 的总和中出现两次，这解释了系数 \(2\)。因此 \(\operatorname{Scal}/(n(n-1))\) 是这 \(n(n-1)/2\) 个截面曲率的平均；其数值与所选正交基无关。常截面曲率 \(k\) 的代数模型满足
\[
	\operatorname{Ric}=(n-1)kg,\qquad
	\operatorname{Scal}=n(n-1)k.
\]

\ProseParagraphBreak

也可以对所有单位径向方向作平均。记 \(\mathrm{d}\sigma\) 为 \(\mathrm{T}_pM\) 的单位球面上的面积测度。由坐标反射和坐标置换的对称性，\(\int v^iv^j\,\mathrm{d}\sigma=0\) 当 \(i\ne j\)，而各 \(\int(v^i)^2\,\mathrm{d}\sigma\) 相同；再用 \(\sum_i(v^i)^2=1\)，得到
\[
	\frac1{\operatorname{Area}(S^{n-1})}
	\int_{S^{n-1}}\operatorname{Ric}_p(v,v)\,\mathrm{d}\sigma(v)
	=\frac1n\operatorname{Scal}(p).
\]
这正是所有径向的体积偏差合成小球体积偏差时出现标量曲率的原因。\cref{g3:thm:small-ball-volume}将在下一节证明：对充分小的 \(r>0\)，
\begin{equation}
	\label{g3:eq:scalar-small-ball-preview}
	\operatorname{Vol}_g(B_d(p,r))
	=\omega_n r^n\left(1-\frac{\operatorname{Scal}(p)}{6(n+2)}r^2+O(r^4)\right),
\end{equation}
其中 \(\omega_n\) 是单位 Euclid \(n\) 球的体积。定向不参与这一公式，体积始终按密度积分。

\ProseParagraphBreak

曲率的这些收缩还满足微分约束。第二 Bianchi 恒等式将曲率在不同方向上的导数联系起来；取迹后，Ricci 的散度恰好由标量曲率的微分控制。

\begin{proposition}[缩并 Bianchi 恒等式]
	\label{g3:prop:contracted-bianchi}
	在光滑黎曼流形 \((M,g)\) 上，定义对称二阶张量 \(S\) 的散度为一形式
	\[
		(\operatorname{div}S)(Y)=\sum_i(\nabla_{e_i}S)(e_i,Y),
	\]
	其中 \(e_i\) 是任意正交单位基。则
	\begin{equation}
		\label{g3:eq:contracted-bianchi}
		\operatorname{div}\operatorname{Ric}=\tfrac12\mathrm{d}\operatorname{Scal}.
	\end{equation}
\end{proposition}
\begin{proof}
	在任意一点取正交法坐标，使该点处 \(g_{ij}=\delta_{ij}\)、\(\nabla_{e_i}e_j=0\)。因 \(\nabla g=0\)，协变微分与用度量作收缩可交换。\eqref{g3:eq:second-bianchi-components}给出
	\[
		\nabla_a\mathrm{Rm}_{ijkl}
		+\nabla_i\mathrm{Rm}_{jakl}
		+\nabla_j\mathrm{Rm}_{aikl}=0.
	\]
	先令 \(k=a\) 并求和。由第一对指标的反对称性及 Ricci 定义，得到
	\[
		\sum_a\nabla_a\mathrm{Rm}_{ijal}
		-\nabla_i\operatorname{Ric}_{jl}
		+\nabla_j\operatorname{Ric}_{il}=0.
	\]
	再令 \(l=i\) 并对 \(i\) 求和。由于 \(\sum_i\mathrm{Rm}_{ijai}=-\operatorname{Ric}_{ja}\)，上式变成
	\[
		-\sum_a\nabla_a\operatorname{Ric}_{ja}
		-\sum_i\nabla_i\operatorname{Ric}_{ji}
		+\nabla_j\operatorname{Scal}=0.
	\]
	Ricci 对称性使前两项相同，所以 \(2(\operatorname{div}\operatorname{Ric})_j=\partial_j\operatorname{Scal}\)。所选点任意，且两边都是一形式，故所得等式在全流形成立。
\end{proof}

例如 \(\operatorname{div}(\operatorname{Ric}-\tfrac12\operatorname{Scal}\,g)=0\)。若某类曲率条件使 Ricci 与度量成比例，这个微分恒等式就会限制比例因子。Schur 定理给出一个基本应用。\footnote{Friedrich Heinrich Schur（1856--1932），德国数学家，研究微分几何、变换群与几何基础；这里的 Schur 定理与群表示论中 Issai Schur 的引理由不同数学家得到。}

\begin{theorem}[Schur]
	\label{g3:thm:schur}
	设 \((M^n,g)\) 为连通光滑黎曼流形，\(n\geq3\)。若在每一点 \(p\)，所有二维子空间的截面曲率具有同一个值 \(k(p)\)，则 \(k\) 在 \(M\) 上为常数。因此 \(M\) 具有常截面曲率。
\end{theorem}
\begin{proof}
	由\cref{g3:cor:pointwise-isotropic-curvature}，在每一点都有\eqref{g3:eq:constant-sectional-algebraic}，从而
	\[
		\operatorname{Ric}=(n-1)kg,\qquad
		\operatorname{Scal}=n(n-1)k.
	\]
	第二式先说明 \(k=\operatorname{Scal}/(n(n-1))\) 是光滑函数。利用 \(\nabla g=0\)，第一式给出 \(\operatorname{div}\operatorname{Ric}=(n-1)\mathrm{d}k\)。代入\eqref{g3:eq:contracted-bianchi}，得
	\[
		(n-1)\mathrm{d}k=\tfrac12n(n-1)\mathrm{d}k,
		\qquad (n-1)(n-2)\mathrm{d}k=0.
	\]
	维数条件使 \(\mathrm{d}k=0\)，所以 \(k\) 局部为常数；连通性使这些局部常数在整个 \(M\) 上相同。
\end{proof}

维数二时，每个切空间本来就只有一个二维子空间，定理的逐点条件不施加额外限制。例如在 \(\mathbb R^2\) 上取
\[
	g=\mathrm{d}r^2+e^{2r^2}\mathrm{d}\theta^2.
\]
由\eqref{g3:eq:warped-surface-curvature}，令 \(f(r)=e^{r^2}\) 后得到
\[
	K(r,\theta)=-\frac{f''(r)}{f(r)}=-2-4r^2,
\]
它不是常数。在这个维数，\(\operatorname{Ric}=Kg\)、\(\operatorname{Scal}=2K\)，缩并 Bianchi 等式两边都等于 \(\mathrm{d}K\)，不再约束 \(K\) 的变化。

\subsection{高维不同正性条件的严格区别}
\label{g3:subsec:0302:04}

曲率算子非负，是说 \(\langle\mathcal R\omega,\omega\rangle\geq0\) 对所有二向量 \(\omega\) 成立；截面曲率非负只在 \(\omega=X\wedge Y\) 这类可分解二向量上作同样的检验。Ricci 非负是二次型 \(\operatorname{Ric}(v,v)\geq0\)，标量曲率非负则只检验一个实数。正曲率算子、正 Ricci 曲率分别指相应二次型正定；正截面曲率要求每个二维截面的值严格为正。

\begin{proposition}
	\label{g3:prop:curvature-positivity-implications}
	在维数 \(n\geq2\) 的黎曼流形的每一点，都有
	\begin{equation}
		\label{g3:eq:curvature-positivity-chain}
		\mathcal R\geq0\ \Longrightarrow\ K\geq0
		\ \Longrightarrow\ \operatorname{Ric}\geq0
		\ \Longrightarrow\ \operatorname{Scal}\geq0.
	\end{equation}
	将其中全部非负号换为严格正号，同样成立。维数二时，四种非负性彼此等价，四种严格正性也彼此等价；维数三时，曲率算子与截面曲率具有相同的非负性和严格正性。
\end{proposition}
\begin{proof}
	第一个箭头来自\eqref{g3:eq:sectional-bivector}，第二个来自\eqref{g3:eq:ricci-sectional-sum}，最后一个来自正交基下取迹。严格正性时，单位方向的 Ricci 和包含 \(n-1\geq1\) 个正数，标量曲率又是 \(n\) 个正数之和。

	维数二时，\(\Lambda^2\mathrm{T}_pM\) 是一维的，曲率算子就是乘以 \(K\)；同时 \(\operatorname{Ric}=Kg\)、\(\operatorname{Scal}=2K\)，故各条件等价。维数三时，每个非零二向量都可分解。具体说，对正交基 \(e_1,e_2,e_3\)，若
	\[
		\omega=a\,e_1\wedge e_2+b\,e_1\wedge e_3+c\,e_2\wedge e_3,
		\qquad a\ne0,
	\]
	则
	\[
		\omega=\left(e_1-\frac ca e_3\right)\wedge(ae_2+be_3).
	\]
	若 \(a=0\)，可重新排列基，使某个非零系数占据该位置。因此在三维，截面曲率的检验已经遍历全部非零二向量，得到所述等价性。
\end{proof}

截面曲率决定整个张量，并不意味着截面上的正性自动成为整个二向量空间上的正性。\cref{g3:prop:sectional-determines-curvature}的恢复过程包含极化求差，求差不会保持正号。下面先在单个内积空间上展示这些区别，再由下一节的局部实现将它们用于度量。

\ProseParagraphBreak

构造代数例子时需要核对第一 Bianchi 恒等式。任取 \(\Lambda^2V\) 上的自伴算子 \(A\)，令 \(T_{ijkl}=\langle A(e_i\wedge e_j),e_k\wedge e_l\rangle\)。它自动具有两对指标的反对称性与交换两对的对称性。第一 Bianchi 恒等式等价于
\begin{equation}
	\label{g3:eq:operator-bianchi-test}
	T_{ijkl}-T_{ikjl}+T_{iljk}=0.
\end{equation}
左侧交换任意相邻两个指标都会变号，这可直接由上述三种对称性核对；例如交换 \(i,j\) 后得到 \(-T_{ijkl}-T_{iljk}+T_{ikjl}\)。所以有重复指标时它自动为零。三维没有四个互异的指标，任意自伴算子因而都通过这个检验；四维则只剩下一个独立条件 \(T_{1234}-T_{1324}+T_{1423}=0\)。

\begin{example}
	\label{g3:exm:ricci-positive-sectional-negative}
	在 \(\mathbb R^3\) 的标准正交基下，令曲率算子在 \(e_1\wedge e_2,e_1\wedge e_3,e_2\wedge e_3\) 上为对角算子，相应特征值记为 \(\kappa_{12},\kappa_{13},\kappa_{23}\)。上一段已经验证它给出代数曲率张量。直接收缩可得
	\[
		\operatorname{Ric}
		=\operatorname{diag}(\kappa_{12}+\kappa_{13},
		\kappa_{12}+\kappa_{23},\kappa_{13}+\kappa_{23}),
	\]
	且 \(\operatorname{Scal}=2(\kappa_{12}+\kappa_{13}+\kappa_{23})\)。取
	\[
		(\kappa_{12},\kappa_{13},\kappa_{23})=(-1,2,2),
	\]
	则 \(\operatorname{Ric}=\operatorname{diag}(1,1,4)\) 正定，但 \(K(e_1,e_2)=-1\)。因此 Ricci 严格为正不能推出截面曲率非负。
\end{example}

\begin{example}
	\label{g3:exm:scalar-positive-ricci-negative}
	在同一三维构造中取
	\[
		(\kappa_{12},\kappa_{13},\kappa_{23})=(-2,3,0).
	\]
	则 \(\operatorname{Ric}=\operatorname{diag}(1,-2,3)\)，而 \(\operatorname{Scal}=2>0\)。标量曲率严格为正仍允许 Ricci 张量有负方向。它记录总和的符号，不能判断每一项的符号。
\end{example}

\begin{example}
	\label{g3:exm:sectional-positive-operator-negative}
	在定向 Euclid 空间 \(\mathbb R^4\) 中，记 \(e_{ij}=e_i\wedge e_j\)，并取下列六个正交单位二向量：
	\[
		\begin{aligned}
			\sigma_1&=(e_{12}+e_{34})/\sqrt2,&
			\tau_1&=(e_{12}-e_{34})/\sqrt2,\\
			\sigma_2&=(e_{13}-e_{24})/\sqrt2,&
			\tau_2&=(e_{13}+e_{24})/\sqrt2,\\
			\sigma_3&=(e_{14}+e_{23})/\sqrt2,&
			\tau_3&=(e_{14}-e_{23})/\sqrt2.
		\end{aligned}
	\]
	定义自伴算子 \(\mathcal R\)，使 \(\sigma_1,\sigma_2,\sigma_3\) 的特征值分别为 \(-1,3,3\)，而三个 \(\tau_i\) 的特征值都为 \(5/3\)。在原二向量基中，相应三个交叉分量为
	\[
		T_{1234}=-\tfrac43,\qquad
		T_{1324}=-\tfrac23,\qquad
		T_{1423}=\tfrac23.
	\]
	它们满足\eqref{g3:eq:operator-bianchi-test}，所以这确实是一个代数曲率张量的算子。

	现在检查所有截面，而不只检查六个坐标截面。记 \(\nu=e_1\wedge e_2\wedge e_3\wedge e_4\)。直接作外积得到
	\[
		\sigma_i\wedge\sigma_j=\delta_{ij}\nu,\qquad
		\tau_i\wedge\tau_j=-\delta_{ij}\nu,\qquad
		\sigma_i\wedge\tau_j=0.
	\]
	若 \(\omega=X\wedge Y\) 是单位可分解二向量，写成 \(\omega=\sum_i a_i\sigma_i+\sum_i b_i\tau_i\)。因为 \(\omega\wedge\omega=0\) 且 \(|\omega|=1\)，有
	\[
		\sum_i a_i^2-\sum_i b_i^2=0,\qquad
		\sum_i a_i^2+\sum_i b_i^2=1.
	\]
	两组平方和因此都为 \(1/2\)，从而
	\[
		\langle\mathcal R\omega,\omega\rangle
		\geq-\sum_i a_i^2+\tfrac53\sum_i b_i^2
		=\tfrac13>0.
	\]
	所有二维截面曲率都严格为正；然而 \(\langle\mathcal R\sigma_1,\sigma_1\rangle=-1\)，所以曲率算子具有负特征值。这里 \(\sigma_1\wedge\sigma_1=\nu\ne0\)，它恰好不是某个二维平面的可分解二向量。
\end{example}

这三个例子首先是单个切空间中的代数数据。\cref{g3:prop:curvature-local-realization}将在下一节证明：任意这样的代数曲率张量都能实现为一个光滑黎曼度量在原点的曲率。上述严格正性与所展示的负方向在充分小的邻域内还会保持。对 Ricci 张量和曲率算子，这来自矩阵分量的连续性；对四维例子的截面曲率，取光滑正交标架后，单位可分解二向量构成固定的紧集，原点处统一的 \(1/3\) 下界可在邻近点保留为某个正下界。因此这些数据给出实际的小开球上的度量反例，说明\eqref{g3:eq:curvature-positivity-chain}的逆向推理在相应维数确实失败；附加紧性或完备性后的整体存在问题需要另外的构造。

\ProseParagraphBreak

相反，若一点处 \(K\geq0\) 且 \(\operatorname{Scal}=0\)，则任一正交基下的非负截面曲率之和为零。将任意指定的二维平面补入一组正交基，就知该平面的曲率也为零；\cref{g3:prop:sectional-determines-curvature}遂给出该点处 \(\mathrm{Rm}=0\)。这个结论同时使用了非负性与零总和。单独的标量曲率为零仍允许不同方向的正负曲率相互抵消。

% 纲要条目排为编号小节标题，并进入目录与书签；说明段仍保留于纲要注释中。

% !TeX root = ../../Book3.tex
% 纲要标题：### 3.3 法坐标中的二阶几何
% 纲要位置：几何学计划纲要.md，第 2210 行。
% 正文状态：初稿；尚未正式审核。

% 纲要内容开始
% 
% * 联络的一阶展开
% * 度量的二阶展开
% * 体积密度和小球体积
% * 展开余项及正则性
% 
% 纲要内容结束

\section{法坐标中的二阶几何}
\label{g3:sec:0303}

在法坐标中心，度量矩阵是单位矩阵，一阶偏导也全部为零。因而仅看一点处的长度和一阶变化，还分辨不出它附近是否弯曲。现在要找的是二阶变化中无法通过这种坐标选择消去的部分。曲率将决定度量的二次项；取矩阵行列式后，这些方向性信息收缩为 Ricci 曲率，再对整个小球积分，就出现标量曲率。

\ProseParagraphBreak

本节的度量始终光滑。固定无边界黎曼流形 \((M^n,g)\) 中一点 \(p\)，在 \(\mathrm T_pM\) 选正交单位基 \((E_1,\ldots,E_n)\)，并用
\[
	q=\exp_p(x^iE_i)
\]
定义足够小邻域内的法坐标。记 \(|x|^2=\sum_i(x^i)^2\)，原点对应 \(p\)。曲率分量沿用本章约定
\[
	R(\partial_i,\partial_j)\partial_b
	=R^a{}_{b i j}\partial_a,\qquad
	R_{ijkl}=g(R(\partial_i,\partial_j)\partial_l,\partial_k).
\]
在原点，因 \(g_{ij}=\delta_{ij}\)，有 \(R^a{}_{b i j}=R_{ijab}\)。下文展开式右侧未注明位置的曲率分量均在 \(p\) 取值；\(O(|x|^m)\) 表示各分量在原点附近的绝对值不超过某个常数乘以 \(|x|^m\)。

\subsection{联络的一阶展开}
\label{g3:subsec:0303:01}

\cref{g3:prop:normal-coordinate-first-order}已经给出 \(\Gamma^a_{ij}(0)=0\)。要确定其一阶项，需要同时使用两类关系：曲率给出联络偏导的反对称部分，径向直线是测地线则给出其完全对称部分。两者合起来恰好确定所有一阶系数。

\begin{proposition}\label{g3:prop:normal-connection-expansion}
	设 \((M^n,g)\) 为光滑无边界黎曼流形，\(x\) 为在 \(p\) 处由正交单位基产生的法坐标。则
	\begin{equation}\label{g3:eq:normal-connection-derivative}
		\partial_k\Gamma^a_{ij}(0)
		=-\frac13\bigl(R^a{}_{i j k}+R^a{}_{j i k}\bigr),
	\end{equation}
	因而
	\begin{equation}\label{g3:eq:normal-connection-expansion}
		\Gamma^a_{ij}(x)
		=-\frac13\bigl(R^a{}_{i j k}+R^a{}_{j i k}\bigr)x^k
		+O(|x|^2).
	\end{equation}
\end{proposition}
\begin{proof}
	在法坐标中，\(t\mapsto tx\) 表示径向测地线。其坐标加速度为零，在时间一代入测地线方程，得到原点附近的恒等式
	\[
		\Gamma^a_{ij}(x)x^ix^j=0.
	\]
	令 \(C^a{}_{ij k}=\partial_k\Gamma^a_{ij}(0)\)。无挠性给出 \(C^a{}_{ij k}=C^a{}_{ji k}\)。把 \(\Gamma\) 的一阶 Taylor 展开代入上式，比较三次齐次项，得到
	\[
		C^a{}_{ij k}x^ix^jx^k=0\qquad\text{对所有 }x.
	\]
	对这个三次多项式取三次偏导，或对其作极化，得
	\begin{equation}\label{g3:eq:normal-connection-symmetric-part}
		C^a{}_{ij k}+C^a{}_{jk i}+C^a{}_{ki j}=0.
	\end{equation}
	这里三项已经包含了完全对称化，因为前两个下指标对称。

	另一方面，原点处 \(\Gamma=0\)，所以\eqref{g3:eq:curvature-coordinate-components}中的二次联络项消失，留下
	\[
		R^a{}_{b i j}=C^a{}_{jb i}-C^a{}_{ib j}.
	\]
	特别地，
	\[
		\begin{aligned}
			R^a{}_{i j k}+R^a{}_{j i k}
			&=C^a{}_{ki j}-C^a{}_{ji k}
			+C^a{}_{kj i}-C^a{}_{ij k}\\
			&=C^a{}_{ki j}+C^a{}_{jk i}-2C^a{}_{ij k}
			=-3C^a{}_{ij k}.
		\end{aligned}
	\]
	最后一步使用\eqref{g3:eq:normal-connection-symmetric-part}。这证明了偏导公式，再对光滑函数 \(\Gamma^a_{ij}\) 使用一阶 Taylor 公式，得到所述展开。
\end{proof}

联络的一阶系数本身依赖坐标，但法坐标的径向条件已经固定了其中可以自由改变的部分。因此在这种坐标中，一阶联络变化可以全部由原点的曲率写出。

\subsection{度量的二阶展开}
\label{g3:subsec:0303:02}

联络与度量之间的联系是度量相容性。它将度量的一阶偏导写成联络系数；再求一次导数，刚得到的一阶联络公式便转化为二阶度量公式。

\begin{theorem}\label{g3:thm:normal-metric-expansion}
	设 \((M^n,g)\) 为光滑无边界黎曼流形，\(x\) 为在 \(p\) 处由正交单位基产生的法坐标。则
	\begin{equation}\label{g3:eq:normal-metric-second-derivative}
		\partial_k\partial_lg_{ij}(0)
		=-\frac13\bigl(R_{ikjl}+R_{iljk}\bigr),
	\end{equation}
	并且
	\begin{equation}\label{g3:eq:normal-metric-expansion}
		g_{ij}(x)=\delta_{ij}-\frac13R_{ikjl}x^kx^l+O(|x|^3).
	\end{equation}
\end{theorem}
\begin{proof}
	度量相容性给出
	\[
		\partial_lg_{ij}=\Gamma^a_{li}g_{aj}+\Gamma^a_{lj}g_{ia}.
	\]
	再对 \(x^k\) 求导，并利用原点处 \(\Gamma=0\)、\(g_{ij}=\delta_{ij}\)，得到
	\[
		\partial_k\partial_lg_{ij}(0)
		=\partial_k\Gamma^j_{li}(0)+\partial_k\Gamma^i_{lj}(0).
	\]
	代入\eqref{g3:eq:normal-connection-derivative}，再将 \(R^a{}_{b i j}\) 改写成四阶分量，右侧变为
	\[
		-\frac13\bigl(R_{ikjl}+R_{lkji}+R_{jkil}+R_{lkij}\bigr).
	\]
	中间的 \(R_{lkji}\) 与最后的 \(R_{lkij}\) 因末两个指标反对称而抵消。成对交换对称性又给出 \(R_{jkil}=R_{iljk}\)，所以得到\eqref{g3:eq:normal-metric-second-derivative}。

	度量的一阶偏导在原点为零，Taylor 公式于是给出
	\[
		g_{ij}(x)=\delta_{ij}
		-\frac16\bigl(R_{ikjl}+R_{iljk}\bigr)x^kx^l+O(|x|^3).
	\]
	第二项中的两个求和式在交换哑指标 \(k,l\) 后相同，便得到\eqref{g3:eq:normal-metric-expansion}。
\end{proof}

这个二次项具有直接的长度含义。在 \(p\) 的切空间中把 \(a=(a^i)\)、\(x=(x^i)\) 看成向量，并在坐标点 \(x\) 考察具有相同分量 \(a^i\) 的坐标切向量，则
\[
	g_{ij}(x)a^ia^j
	=|a|^2-\frac13\mathrm{Rm}_p(a,x,a,x)
	+O(|x|^3)|a|^2.
\]
若 \(a\) 与 \(x\) 垂直，二次修正由它们张成平面的截面曲率决定。若 \(a\) 与 \(x\) 平行，这个曲率项因反对称性消失，与 Gauss 引理所给的径向长度保持相符。

\begin{example}\label{g3:exm:constant-curvature-normal-quadratic}
	设 \(n\ge2\)。若 \(p\) 处所有截面曲率等于常数 \(\kappa\)，由\cref{g3:cor:pointwise-isotropic-curvature}有
	\[
		R_{ikjl}=\kappa(\delta_{ij}\delta_{kl}-\delta_{il}\delta_{kj}),
	\]
	故二阶度量展开为
	\[
		g_{ij}(x)=\delta_{ij}
		-\frac{\kappa}{3}\bigl(|x|^2\delta_{ij}-x^ix^j\bigr)+O(|x|^3).
	\]
	对垂直于 \(x\) 的单位向量 \(a\)，其在坐标点 \(x\) 的长度平方为 \(1-\kappa|x|^2/3+O(|x|^3)\)。正曲率使这一横向长度在二阶上小于平坦模型，负曲率则使它大于平坦模型。径向与横向表现不同，因此不能把二阶修正理解为对所有方向的统一缩放。
\end{example}

固定一个二维切平面，再把其中的小圆盘经指数映射送到流形，可以直接观测这个截面的曲率。得到的是嵌入二维圆盘及其诱导度量；其边界是参数圆周的像。

\begin{proposition}\label{g3:prop:normal-plane-disk-geometry}
	设 \((M^n,g)\) 为光滑无边界黎曼流形，\(n\ge2\)，\(\sigma\subset\mathrm T_pM\) 为二维线性子空间。对足够小的 \(r>0\)，令
	\[
		D_r=\exp_p\{v\in\sigma:|v|_g\le r\},\qquad
		\partial D_r=\exp_p\{v\in\sigma:|v|_g=r\}.
	\]
	这里 \(\partial D_r\) 表示这个二维圆盘的边界。其诱导面积与边界长度满足
	\begin{equation}\label{g3:eq:normal-plane-disk-geometry}
		\begin{aligned}
			\operatorname{Area}_g(D_r)
			&=\pi r^2\left(1-\frac{K_p(\sigma)}{12}r^2+O(r^4)\right),\\
			\operatorname{Length}_g(\partial D_r)
			&=2\pi r\left(1-\frac{K_p(\sigma)}6r^2+O(r^4)\right).
		\end{aligned}
	\end{equation}
\end{proposition}
\begin{proof}
	选取正交单位基，使 \(\sigma\) 由 \(E_1,E_2\) 张成。令 \(z=(x^1,x^2)\)，其余坐标为零。指数映射在一个较大的切球上为微分同胚，所以限制后得到嵌入圆盘；其诱导度量矩阵为 \(h_{ab}(z)=g_{ab}(z,0)\)，\(1\le a,b\le2\)。记 \(K=K_p(\sigma)\)。曲率限制到这个二维空间后，只有 \(R_{1212}=K\) 一个独立分量，故\eqref{g3:eq:normal-metric-expansion}给出
	\[
		h_{ab}(z)=\delta_{ab}
		-\frac K3\bigl(|z|^2\delta_{ab}-z^az^b\bigr)
		+P_{ab,3}(z)+O(|z|^4),
	\]
	其中 \(P_{ab,3}\) 为光滑系数的三次齐次 Taylor 项。对单位矩阵附近的矩阵 \(I+H\)，有 \(\sqrt{\det(I+H)}=1+\tfrac12\operatorname{tr}H+O(\|H\|^2)\)。这里 \(H=O(|z|^2)\)，所以诱导面积密度为
	\[
		J_\sigma(z)=\sqrt{\det h(z)}
		=1-\frac K6|z|^2+P_3(z)+O(|z|^4),
	\]
	其中 \(P_3\) 仍为三次齐次多项式。在半径 \(r\) 的平面圆盘上积分，奇函数 \(P_3\) 的积分为零，而 \(\int_{|z|\le r}|z|^2\,\mathrm dz=\pi r^4/2\)。四阶余项积分为 \(O(r^6)\)，得到面积公式。

	对边界令 \(u(\theta)=(\cos\theta,\sin\theta)\)、\(v(\theta)=(-\sin\theta,\cos\theta)\)。其坐标速度为 \(rv(\theta)\)，由 \(u,v\) 为正交单位向量，有
	\[
		|\dot c_r(\theta)|_g^2
		=r^2\left(1-\frac K3r^2+O(r^3)\right).
	\]
	为同时控制积分余项，使用 Gauss 引理：圆盘上的径向方向长度为一，并与角向方向正交。因此极坐标中的面积密度既等于 \(|\dot c_r(\theta)|_g\,\mathrm dr\,\mathrm d\theta\)，也等于 \(rJ_\sigma(ru(\theta))\,\mathrm dr\,\mathrm d\theta\)，从而
	\[
		\operatorname{Length}_g(\partial D_r)
		=r\int_0^{2\pi}J_\sigma(ru(\theta))\,\mathrm d\theta.
	\]
	在 \(\theta\mapsto\theta+\pi\) 下，三次项变号，积分为零；余下四阶项贡献 \(O(r^5)\)。代入二次项便得长度公式。
\end{proof}

这些圆盘由同一个点的径向测地线组成；盘内其余方向的内在测地线未必是环境测地线，因此不能据此把圆盘视为全测地子流形。面积和长度是对这个指定二维指数像作出的测量。当 \(n>2\) 时，它也不同于环境中的 \(n\) 维测地球。截面曲率正是这类二维测量的首个非平坦系数。

\ProseParagraphBreak
二阶公式还有一个反向用途。前两节讨论的代数曲率张量是否真能出现在某个度量的一点处，可以通过直接指定度量的二次项来回答。

\begin{proposition}\label{g3:prop:curvature-local-realization}
	设 \(A\) 为 \(n\) 维实 Euclid 空间上的代数曲率张量，即满足\cref{g3:def:algebraic-curvature-tensor}中的两对反对称性、成对交换对称性和第一 Bianchi 恒等式。则存在原点的一个邻域及其上的光滑黎曼度量，使原点处度量为 Euclid 内积，曲率张量为 \(A\)。具体可取
	\begin{equation}\label{g3:eq:curvature-local-realization-metric}
		g_{ij}(x)=\delta_{ij}-\frac13A_{ikjl}x^kx^l
	\end{equation}
	并将定义域限制在足够小的开球内。
\end{proposition}
\begin{proof}
	先核对它确实是度量。由成对交换对称性和哑指标互换，
	\[
		A_{jkil}x^kx^l=A_{iljk}x^kx^l=A_{ikjl}x^kx^l,
	\]
	故 \(g_{ij}=g_{ji}\)。因为 \(A\) 固定，存在常数 \(C\ge0\)，使任意 \(v\in\mathbb R^n\) 满足
	\[
		\left|\frac13A_{ikjl}v^iv^jx^kx^l\right|
		\le C|x|^2|v|^2.
	\]
	在 \(C|x|^2<1/2\) 的小球内，\(g(v,v)\ge|v|^2/2\)，所以正定。系数是多项式，因而光滑。原点处 \(g_{ij}=\delta_{ij}\)、\(\partial_kg_{ij}=0\)，故 \(\Gamma^a_{ij}(0)=0\)，且
	\[
		\partial_k\partial_lg_{ij}(0)
		=-\frac13\bigl(A_{ikjl}+A_{iljk}\bigr).
	\]
	此处的 \(g(0)=I\)、\(\partial g(0)=0\) 正好满足\eqref{g3:eq:curvature-normal-second-derivatives}推导时所用的点处条件。将二阶偏导代入该曲率公式，整理两对反对称性与成对交换对称性，得到
	\[
		R^g_{ijkl}(0)
		=\frac13\bigl(2A_{ijkl}+A_{ikjl}-A_{iljk}\bigr).
	\]
	第一 Bianchi 恒等式为
	\[
		A_{ijkl}+A_{jkil}+A_{kijl}=0.
	\]
	其中 \(A_{jkil}=A_{iljk}\)、\(A_{kijl}=-A_{ikjl}\)，所以 \(A_{ikjl}-A_{iljk}=A_{ijkl}\)。代回即得 \(R^g_{ijkl}(0)=A_{ijkl}\)，完成构造。
\end{proof}

这个构造实现的是指定点的曲率数据。若某个曲率不等式在 \(A\) 处严格成立，则由曲率的连续性，将邻域再缩小后，该严格不等式仍成立；涉及所有单位向量或单位二维截面时，所需的统一余量来自相应单位参数空间的紧性。因此，严格的代数正性反例也能给出局部黎曼度量的反例。

\subsection{体积密度和小球体积}
\label{g3:subsec:0303:03}

度量矩阵决定每个方向的长度，体积密度则把这些方向的伸缩合成为一个标量。在法坐标中，令
\[
	J_p(x)=\sqrt{\det(g_{ij}(x))},\qquad
	\mu_g=J_p(x)|\mathrm dx^1\cdots\mathrm dx^n|.
\]
这里的密度定义不需要流形可定向。

\begin{proposition}\label{g3:prop:normal-volume-density}
	设 \((M^n,g)\) 为光滑无边界黎曼流形，在 \(p\) 处采用正交单位基产生的法坐标。则
	\begin{equation}\label{g3:eq:normal-volume-density}
		J_p(x)=1-\frac16\operatorname{Ric}_{kl}(p)x^kx^l+O(|x|^3).
	\end{equation}
\end{proposition}
\begin{proof}
	写 \(G=I+H\)。由\eqref{g3:eq:normal-metric-expansion}，有
	\[
		H_{ij}=-\frac13R_{ikjl}x^kx^l+O(|x|^3),
		\qquad H=O(|x|^2).
	\]
	行列式是矩阵元的多项式，故在单位矩阵附近，
	\[
		\det(I+H)=1+\operatorname{tr}H+O(\|H\|^2),
		\qquad
		\sqrt{\det(I+H)}=1+\tfrac12\operatorname{tr}H+O(\|H\|^2).
	\]
	另一方面，按本章的 Ricci 收缩约定，
	\[
		\operatorname{tr}H
		=-\frac13\sum_iR_{ikil}x^kx^l+O(|x|^3)
		=-\frac13\operatorname{Ric}_{kl}(p)x^kx^l+O(|x|^3).
	\]
	因为 \(\|H\|^2=O(|x|^4)\)，代回便得所述展开。
\end{proof}

沿固定单位方向 \(u\) 取 \(x=ru\)，则密度的二阶项为 \(-\operatorname{Ric}_p(u,u)r^2/6\)。这说明 Ricci 曲率测量的是沿该方向向外展开时，横向体积相对于 Euclid 模型的二阶变化。将所有方向平均，才得到标量曲率的作用。

\begin{theorem}\label{g3:thm:small-ball-volume}
	设 \((M^n,g)\) 为光滑无边界黎曼流形，\(p\in M\)，\(d\) 为 \(p\) 所在连通分支的内在距离，\(\omega_n\) 表示 \(\mathbb R^n\) 中单位球的 Euclid 体积。则当 \(r\downarrow0\) 时，内在距离球的体积满足
	\begin{equation}\label{g3:eq:small-ball-volume}
		\operatorname{Vol}_g(B_d(p,r))
		=\omega_n r^n\left(1-\frac{\operatorname{Scal}(p)}{6(n+2)}r^2+O(r^4)\right).
	\end{equation}
	若 \(n\ge2\)，足够小的距离球边界是光滑超曲面，其面积满足
	\begin{equation}\label{g3:eq:small-sphere-area}
		\operatorname{Area}_g(\partial B_d(p,r))
		=n\omega_n r^{n-1}\left(1-\frac{\operatorname{Scal}(p)}{6n}r^2+O(r^4)\right).
	\end{equation}
\end{theorem}
\begin{proof}
	由\cref{g3:lem:uniform-normal-neighborhood}，对足够小的 \(r\)，\(B_d(p,r)=\exp_p B_p(0,r)\)，且指数映射在更大的切球上为微分同胚。因此体积可以在普通 Euclid 球上积分：
	\[
		\operatorname{Vol}_g(B_d(p,r))=\int_{|x|<r}J_p(x)\,\mathrm dx.
	\]
	为了确定积分后的余项阶数，需把密度再展开一阶。\(J_p\) 光滑，一次项为零，二次项已经由\cref{g3:prop:normal-volume-density}确定，所以存在三次齐次多项式 \(P_{3,p}\)，使
	\begin{equation}\label{g3:eq:normal-density-cubic-remainder}
		J_p(x)=1-\frac16\operatorname{Ric}_{kl}(p)x^kx^l
		+P_{3,p}(x)+O(|x|^4).
	\end{equation}
	三次齐次项满足 \(P_{3,p}(-x)=-P_{3,p}(x)\)，在中心对称的球上积分为零。对二次项，由反射和坐标置换对称性，非对角积分为零，各对角积分相等；再把它们相加，得到
	\[
		\int_{|x|<r}x^kx^l\,\mathrm dx
		=\delta_{kl}\frac{\omega_n}{n+2}r^{n+2}.
	\]
	故二次项的积分为 \(-\omega_n\operatorname{Scal}(p)r^{n+2}/(6(n+2))\)。四阶余项的绝对值积分至多为常数乘以 \(r^{n+4}\)，除去主项 \(\omega_n r^n\) 后即得\eqref{g3:eq:small-ball-volume}。

	再求球面面积。令 \(x=ru\)、\(u\in S^{n-1}\)，并以 \(\mathrm d\sigma(u)\) 表示单位 Euclid 球面的面积密度。Gauss 引理给出径向单位向量与距离球面正交，因而在这些极坐标中，体积密度等于 \(\mathrm dr\) 与各距离球面面积密度的乘积。同时 Euclid 极坐标变换给出 \(\mathrm dx=r^{n-1}\mathrm dr\,\mathrm d\sigma\)，所以
	\[
		\operatorname{Area}_g(\partial B_d(p,r))
		=r^{n-1}\int_{S^{n-1}}J_p(ru)\,\mathrm d\sigma(u).
	\]
	将\eqref{g3:eq:normal-density-cubic-remainder}代入，三次项仍因 \(u\mapsto-u\) 对称而积分为零；又有
	\[
		\int_{S^{n-1}}1\,\mathrm d\sigma=n\omega_n,
		\qquad
		\int_{S^{n-1}}u^ku^l\,\mathrm d\sigma=\delta_{kl}\omega_n.
	\]
	这便给出\eqref{g3:eq:small-sphere-area}。面积公式直接来自极坐标密度积分，所需的余项估计也在该积分中取得。
\end{proof}

若 \(\operatorname{Scal}(p)>0\)，足够小的测地球体积小于同半径 Euclid 球；若 \(\operatorname{Scal}(p)<0\)，则大于它。标量曲率为零时，二次修正消失，体积差异由更高阶项决定。这个平均结果并不决定每个方向的 Ricci 曲率符号。

\begin{example}\label{g3:exm:surface-small-ball-check}
	二维时，\(\operatorname{Scal}=2K\)、\(\omega_2=\pi\)，所以
	\[
		\begin{aligned}
			\operatorname{Area}_g(B_d(p,r))
			&=\pi r^2\left(1-\frac{K(p)}{12}r^2+O(r^4)\right),\\
			\operatorname{Length}_g(\partial B_d(p,r))
			&=2\pi r\left(1-\frac{K(p)}6r^2+O(r^4)\right).
		\end{aligned}
	\]
	可以用刚构造的二次度量作具体核验。对实常数 \(\kappa\)，在原点附近令
	\[
		g_{ij}(x)=\delta_{ij}-\frac{\kappa}{3}
		\bigl(|x|^2\delta_{ij}-x^ix^j\bigr).
	\]
	它是\cref{g3:prop:curvature-local-realization}中的二维模型，故 \(K(0)=\kappa\)。改写成极坐标得
	\[
		g=\mathrm dr^2+r^2\left(1-\frac{\kappa r^2}{3}\right)\mathrm d\theta^2.
	\]
	在足够小的范围内括号内为正。翘曲积联络公式说明径向曲线以 \(r\) 为单位速参数是测地线；原来的直角坐标系数在原点光滑，故这些径向线与原点的测地线初值相接，\((x^1,x^2)\) 就是原点的法坐标。此时
	\[
		J_0(x)=\sqrt{1-\frac{\kappa|x|^2}{3}}
		=1-\frac{\kappa}{6}|x|^2-\frac{\kappa^2}{72}|x|^4+O(|x|^6).
	\]
	沿半径积分或沿圆周积分，分别得到
	\[
		\begin{aligned}
			\operatorname{Area}_g(B_d(0,r))
			&=\pi r^2\left(1-\frac{\kappa}{12}r^2
			-\frac{\kappa^2}{216}r^4+O(r^6)\right),\\
			\operatorname{Length}_g(\partial B_d(0,r))
			&=2\pi r\left(1-\frac{\kappa}{6}r^2
			-\frac{\kappa^2}{72}r^4+O(r^6)\right).
		\end{aligned}
	\]
	两式的二阶相对修正只用 \(K(0)=\kappa\) 即已确定。球的面积与边界长度采用不同系数，因为前者还对径向变量作了一次积分；再高阶的系数则取决于这个具体度量。
\end{example}

\subsection{展开余项及正则性}
\label{g3:subsec:0303:04}

上述公式同时记录了系数和误差，两者承担不同职责。曲率决定二阶系数；误差界还依赖附近度量系数的更高阶导数，不能只由 \(R(p)\) 的值控制。

\ProseParagraphBreak

在一个固定法坐标图中，令
\[
	E_{ij}(x)=g_{ij}(x)-\delta_{ij}+\frac13R_{ikjl}(p)x^kx^l.
\]
它在原点的零阶、一阶和二阶导数全部为零。对每个满足 \(|\alpha|\le2\) 的多重指标，分别对 \(\partial^\alpha E_{ij}\) 使用 Taylor 公式，得到
\begin{equation}\label{g3:eq:normal-metric-remainder-derivatives}
	\partial^\alpha E_{ij}(x)=O(|x|^{3-|\alpha|}).
\end{equation}
这些导数估计来自光滑函数的 Taylor 余项，而非对一个任意 \(O\) 记号形式求导。同样，若要在小球体积公式中保留相对 \(O(r^4)\)，就应使用\eqref{g3:eq:normal-density-cubic-remainder}的三次项和四阶余项；仅对密度中的 \(O(|x|^3)\) 作绝对值估计，只能得到相对 \(O(r^3)\)。

\begin{proposition}\label{g3:prop:normal-expansion-uniform-remainders}
	设 \((M,g)\) 为光滑无边界黎曼流形，\(K\subset M\) 紧。则存在共同半径 \(r_0>0\)，使本节的联络、度量、密度、小球体积展开，以及在 \(n\ge2\) 时的小球面面积展开，对所有中心 \(p\in K\)、\(0<r<r_0\) 都成立，且各余项中的常数可与 \(p\) 无关。
\end{proposition}
\begin{proof}
	由\cref{g3:lem:uniform-normal-neighborhood}，\(K\) 上存在统一可用的法坐标半径。取有限个较小的坐标邻域覆盖 \(K\)，使各自闭包紧含于一个有光滑正交单位标架 \((E_1(p),\ldots,E_n(p))\) 的较大邻域。在每个这样的邻域中，
	\[
		F(p,x)=\exp_p(x^iE_i(p))
	\]
	在共同的小 \(x\) 球上光滑依赖 \((p,x)\)。这由\cref{g3:prop:exponential-local}的指数映射联合光滑性和标架光滑性得到。拉回度量系数
	\[
		G_{ij}(p,x)=g_{F(p,x)}
		\bigl(\mathrm d_xF(p,x)e_i,\mathrm d_xF(p,x)e_j\bigr)
	\]
	因此也联合光滑。将 \(x\) 球再缩小一层后，各底邻域的闭包与闭 \(x\) 球的乘积紧，所以证明中使用的有限阶导数均有统一上界，度量正定性也有统一下界。Taylor 公式、行列式及平方根的导数估计便给出统一余项界。

	有限个邻域取最小半径、最大误差常数，即得对 \(K\) 一致的结论。改用同一点处的其他正交单位基，只对 \(x\) 作正交线性变换，故不改变 \(|x|\) 的大小；相应张量范数也保持不变。体积和面积的积分估计由这些统一界直接得到。
\end{proof}

这里始终假设原度量光滑，从而指数映射与法坐标中的度量系数也光滑。若降低原度量的正则性，需要分别核对测地线方程、指数映射和换坐标后的系数具有多少阶导数。就 Taylor 估计本身而言，已给定坐标中的 \(C^2\) 系数只能直接保证二阶展开的 \(o(|x|^2)\) 余项；三阶导数有界给出 \(O(|x|^3)\)，而密度展开至三次项后控制 \(O(|x|^4)\) 需要相应的四阶导数界。原坐标中的可微阶数不能未经检查就沿用到法坐标中。

% !TeX root = ../../Book3.tex
% 纲要标题：### 3.4 模型与代数分解
% 纲要位置：几何学计划纲要.md，第 2114 行。
% 正文状态：未撰写；本文件目前是框架。

% 纲要内容开始
% 
% * 空间形式、乘积和翘曲积
% * 共形变化和标量曲率
% * Weyl 张量及低维特例
% * 由模型检验符号与等号
% 
% ---
% 
% 纲要内容结束

\section{模型与代数分解}
\label{g3:sec:0304}

待撰写

\subsection{空间形式、乘积和翘曲积}
\label{g3:subsec:0304:01}
待撰写

\subsection{共形变化和标量曲率}
\label{g3:subsec:0304:02}
待撰写

\subsection{Weyl 张量及低维特例}
\label{g3:subsec:0304:03}
待撰写

\subsection{由模型检验符号与等号}
\label{g3:subsec:0304:04}
待撰写
% 纲要条目排为编号小节标题，并进入目录与书签；说明段仍保留于纲要注释中。


\begin{chaptersummary}
	曲率是协变求导的交换差扣除 Lie 括号方向导数后留下的张量。度量相容性给出反对称性，无挠性给出第一 Bianchi 恒等式；第二 Bianchi 恒等式进一步限制曲率的变化，其收缩产生 Ricci 散度与标量曲率微分之间的关系。把四阶张量视为二向量空间上的自伴算子后，截面曲率只检测可分解二向量；这一点解释了高维中截面正性与算子正性之间的差别。

	\ProseParagraphBreak

	法坐标的径向条件与曲率公式共同确定联络的一阶项，度量相容性再确定度量的二阶项。二维切平面指数像的周长和面积感受到相应的截面曲率，法坐标体积密度的二次项由 Ricci 曲率控制，小球体积的方向平均则由标量曲率控制。Taylor 余项与对称积分使这些解释成为带有误差阶数的局部公式；紧中心集上的共同法坐标范围还给出一致估计。

	\ProseParagraphBreak

	乘积与翘曲积使曲率的不同来源可以逐项计算，共形变化把长度尺度的导数转化为曲率修正。Weyl 分解将曲率拆成相互正交的三部分：标量部分、无迹 Ricci 部分和无迹四阶部分。二维曲率由一个数决定，三维由 Ricci 决定，四维起则可以有 Ricci 收缩看不见的曲率。标准模型既核验了符号，也说明局部体积或某一次收缩不能单独恢复全部几何。
\end{chaptersummary}

\BookExercises


\begin{exercise}\label{g3:ex:constant-metric-curvature-scaling}
	设 \((M^n,g)\) 为光滑黎曼流形，\(a>0\) 为常数，\(\widetilde g=a^2g\)。将两种度量下的切空间、张量空间和二向量空间按恒等映射识别。
	\begin{enumerate}
		\item 求 \((1,3)\) 型曲率 \(R\)、四阶协变曲率 \(\mathrm{Rm}\)、协变 Ricci 张量及标量曲率在此次缩放下的变化。
		\item 求曲率算子 \(\mathcal R\) 与曲率平方范数 \(|\mathrm{Rm}|^2\) 的变化。说明在固定基与分别单位化的基中，曲率分量的比较为何不同。
	\end{enumerate}
\end{exercise}
\begin{solution}
	1. 原联络 \(\nabla\) 仍然无挠，而且
	\[
		X\widetilde g(Y,Z)
		=a^2Xg(Y,Z)
		=\widetilde g(\nabla_XY,Z)+\widetilde g(Y,\nabla_XZ).
	\]
	由 \cref{g3:thm:levi-civita} 的唯一性，\(\widetilde\nabla=\nabla\)。曲率作为由联络定义的 \((1,3)\) 型张量因而不变：
	\[
		\widetilde R(X,Y)Z=R(X,Y)Z.
	\]
	四阶张量还含一次度量配对，故
	\[
		\widetilde{\mathrm{Rm}}(X,Y,Z,W)
		=\widetilde g(\widetilde R(X,Y)W,Z)
		=a^2\mathrm{Rm}(X,Y,Z,W).
	\]
	在同一坐标基中，\(\widetilde g^{ik}=a^{-2}g^{ik}\)，所以
	\[
		\widetilde{\operatorname{Ric}}_{jl}
		=\widetilde g^{ik}\widetilde R_{ijkl}
		=g^{ik}R_{ijkl}=\operatorname{Ric}_{jl},
		\qquad
		\widetilde{\operatorname{Scal}}
		=a^{-2}\operatorname{Scal}.
	\]
	这里不变的是协变二阶 Ricci 张量；将它用度量升起一个指标后，所得自伴算子的特征值也乘以 \(a^{-2}\)。

	2. 二向量空间上的诱导内积含两次度量，满足
	\[
		\langle\xi,\eta\rangle_{\Lambda^2,\widetilde g}
		=a^4\langle\xi,\eta\rangle_{\Lambda^2,g}.
	\]
	由曲率算子的定义，对可分解二向量 \(\xi=X\wedge Y\)、\(\eta=Z\wedge W\)，
	\[
		a^4\langle\widetilde{\mathcal R}\xi,\eta\rangle_{\Lambda^2,g}
		=a^2\langle\mathcal R\xi,\eta\rangle_{\Lambda^2,g}.
	\]
	这类二向量张成整个空间，内积又非退化，因此
	\[
		\widetilde{\mathcal R}=a^{-2}\mathcal R.
	\]
	在二形式空间上用各自的度量同构定义曲率算子时，结论相同：\(\widetilde\flat_2=a^4\flat_2\)、\(\widetilde\sharp_2=a^{-4}\sharp_2\)，两个同构的尺度在共轭时抵消。

	设 \((e_i)\) 为 \(g\) 正交单位基，则 \((\widetilde e_i)=(a^{-1}e_i)\) 为 \(\widetilde g\) 正交单位基。在同一组固定向量 \(e_i\) 上，\(\widetilde{\mathrm{Rm}}\) 的分量是原来的 \(a^2\) 倍；在各自的单位基上则有
	\[
		\widetilde{\mathrm{Rm}}
			(\widetilde e_i,\widetilde e_j,\widetilde e_k,\widetilde e_l)
		=a^{-2}\mathrm{Rm}(e_i,e_j,e_k,e_l).
	\]
	因此
	\[
		|\widetilde{\mathrm{Rm}}|_{\widetilde g}^2
		=a^{-4}|\mathrm{Rm}|_g^2.
	\]
	同样，每个固定二维平面的截面曲率变为原来的 \(a^{-2}\) 倍。张量本身的缩放与分量比较中基向量的缩放共同决定这些公式。
\end{solution}

\begin{exercise}\label{g3:ex:algebraic-curvature-dimension}
	设 \(V\) 为 \(n\) 维实向量空间，\(n\ge2\)。将 \(\mathcal S=\operatorname{Sym}^2(\Lambda^2V^*)\) 视为满足两对槽位内反对称及成对交换对称的四阶协变张量空间。对 \(A\in\mathcal S\) 定义
	\[
		(\mathfrak bA)(X,Y,Z,W)
		=A(X,Y,Z,W)+A(Y,Z,X,W)+A(Z,X,Y,W).
	\]
	\begin{enumerate}
		\item 证明 \(\mathfrak bA\) 是交替四形式，且线性映射
		\(\mathfrak b:\mathcal S\to\Lambda^4V^*\) 满射。
		\item 证明代数曲率张量空间就是 \(\ker\mathfrak b\)，并求其维数。特别求 \(n=2,3,4\) 时的数值。
	\end{enumerate}
\end{exercise}
\begin{solution}
	1. 利用成对交换与第一对的反对称性，可将定义写成
	\[
		(\mathfrak bA)(X,Y,Z,W)
		=A(X,Y,Z,W)-A(X,Z,Y,W)+A(X,W,Y,Z).
	\]
	记右侧为 \(B(X,Y,Z,W)\)。例如，交换前两个输入并再次使用成对交换，得到
	\[
		\begin{aligned}
			B(Y,X,Z,W)
			&=-A(X,Y,Z,W)-A(Y,Z,X,W)+A(Y,W,X,Z)\\
			&=-A(X,Y,Z,W)-A(X,W,Y,Z)+A(X,Z,Y,W)\\
			&=-B(X,Y,Z,W).
		\end{aligned}
	\]
	其余两个相邻换位也可直接展开：
	\[
		\begin{aligned}
			B(X,Z,Y,W)
			&=A(X,Z,Y,W)-A(X,Y,Z,W)+A(X,W,Z,Y)\\
			&=-B(X,Y,Z,W),\\
			B(X,Y,W,Z)
			&=A(X,Y,W,Z)-A(X,W,Y,Z)+A(X,Z,Y,W)\\
			&=-B(X,Y,Z,W).
		\end{aligned}
	\]
	相邻换位生成全部置换，故 \(B\) 是交替四形式。

	反过来，任取 \(\eta\in\Lambda^4V^*\)，也把它看作四阶张量。交换每一对内的输入使它变号，交换两对输入是一个偶置换，所以 \(\eta\in\mathcal S\)。循环前三个槽位也是偶置换，因此
	\[
		\mathfrak b\eta=3\eta,\qquad
		\mathfrak b\left(\frac13\eta\right)=\eta.
	\]
	这证明满射。当 \(n<4\) 时，目标空间为零，上述结论仍成立。

	2. \(\mathcal S\) 已包含定义中的两对反对称性和成对交换对称性，剩下的第一 Bianchi 恒等式恰为 \(\mathfrak bA=0\)。所以其核正是代数曲率张量空间。

	令 \(N=\dim\Lambda^2V^*=n(n-1)/2\)。在一组二形式基中，\(\mathcal S\) 对应实对称 \(N\times N\) 矩阵，因此维数为 \(N(N+1)/2\)。满射的秩为 \(\binom n4\)，秩与零度关系给出
	\[
		\begin{aligned}
			\dim\ker\mathfrak b
			&=\frac{N(N+1)}2-\binom n4\\
			&=\frac{n(n-1)\bigl(n(n-1)+2\bigr)}8
				-\frac{n(n-1)(n-2)(n-3)}{24}\\
			&=\frac{n^2(n^2-1)}{12}.
		\end{aligned}
	\]
	这里约定 \(n<4\) 时 \(\binom n4=0\)。于是二维、三维、四维的维数依次为
	\[
		1,\qquad6,\qquad20.
	\]
	特别地，三维中每个对称的二向量双线性型都满足第一 Bianchi 恒等式；四维中则恰有一个独立的线性约束。
\end{solution}

\begin{exercise}\label{g3:ex:curvature-covariant-tensors}
	设 \((M,g)\) 为光滑黎曼流形，\(\nabla\) 为其 Levi-Civita 联络，并在协变张量上使用诱导联络
	\[
		(\nabla_XT)(Z_1,\ldots,Z_k)
		=X\bigl(T(Z_1,\ldots,Z_k)\bigr)
		-\sum_{a=1}^kT(Z_1,\ldots,\nabla_XZ_a,\ldots,Z_k).
	\]
	对函数或协变张量场 \(T\)，记
	\[
		\mathcal D_{X,Y}T
		=\nabla_X\nabla_YT-\nabla_Y\nabla_XT-\nabla_{[X,Y]}T.
	\]
	\begin{enumerate}
		\item 证明对任意一形式 \(\alpha\)，
		\[
			(\mathcal D_{X,Y}\alpha)(Z)=-\alpha(R(X,Y)Z).
		\]
		\item 对一般协变 \(k\) 阶张量，证明
		\[
			(\mathcal D_{X,Y}T)(Z_1,\ldots,Z_k)
			=-\sum_{a=1}^k
				T(Z_1,\ldots,R(X,Y)Z_a,\ldots,Z_k).
		\]
		\item 说明函数的协变二阶导数可以交换，但一形式的协变二阶导数一般不能交换，并在正曲率圆球面上给出一个具体的非零值。
	\end{enumerate}
\end{exercise}
\begin{solution}
	1. 把诱导联络的定义使用两次，得到
	\[
		\begin{aligned}
			(\nabla_X\nabla_Y\alpha)(Z)
			&=X\bigl((\nabla_Y\alpha)(Z)\bigr)
				-(\nabla_Y\alpha)(\nabla_XZ)\\
			&=XY\bigl(\alpha(Z)\bigr)
				-X\bigl(\alpha(\nabla_YZ)\bigr)
				-Y\bigl(\alpha(\nabla_XZ)\bigr)
				+\alpha(\nabla_Y\nabla_XZ).
		\end{aligned}
	\]
	交换 \(X,Y\) 后，两个中间项相消。再减去
	\[
		(\nabla_{[X,Y]}\alpha)(Z)
		=[X,Y]\bigl(\alpha(Z)\bigr)-\alpha(\nabla_{[X,Y]}Z),
	\]
	作用在标量函数 \(\alpha(Z)\) 上的三项也相消，余下
	\[
		(\mathcal D_{X,Y}\alpha)(Z)
		=-\alpha\bigl(
			\nabla_X\nabla_YZ-\nabla_Y\nabla_XZ-\nabla_{[X,Y]}Z
			\bigr),
	\]
	即所求公式。负号来自协变张量对输入向量变化的扣除。

	2. 对函数 \(f\)，有
	\[
		\mathcal D_{X,Y}f=XYf-YXf-[X,Y]f=0.
	\]
	诱导联络满足张量积的乘积规则。为检查两次求导后的规则，展开
	\[
		\begin{aligned}
			\nabla_X\nabla_Y(S\otimes T)
			={}&(\nabla_X\nabla_YS)\otimes T
				+(\nabla_YS)\otimes(\nabla_XT)\\
			&+(\nabla_XS)\otimes(\nabla_YT)
				+S\otimes(\nabla_X\nabla_YT).
		\end{aligned}
	\]
	交换 \(X,Y\) 时，两个交叉项相消；再减去括号方向的导数，得到
	\[
		\mathcal D_{X,Y}(S\otimes T)
		=(\mathcal D_{X,Y}S)\otimes T
			+S\otimes(\mathcal D_{X,Y}T).
	\]
	对 \(S=\alpha^1\otimes\cdots\otimes\alpha^k\) 逐次使用此式，再代入第 1 问的一形式公式，恰好得到每个槽位各插入一次 \(-R(X,Y)\) 的和。任意协变张量在局部余标架中都是这种纯张量的光滑函数线性组合；函数上的 \(\mathcal D_{X,Y}\) 为零，故同一公式对一般 \(T\) 成立。

	3. 函数 \(f\) 的协变二阶导数为
	\[
		H_f(X,Y)=X(Yf)-(\nabla_XY)f.
	\]
	利用无挠性，
	\[
		H_f(X,Y)-H_f(Y,X)
		=[X,Y]f-(\nabla_XY-\nabla_YX)f=0.
	\]
	这里需要区分普通迭代方向导数与张量二阶导数：\(XYf-YXf=[X,Y]f\)，对交换的坐标向量场这就是普通混合偏导相等。

	对一形式定义张量二阶导数
	\[
		(\nabla^2\alpha)(X,Y;Z)
		=\bigl(\nabla_X\nabla_Y\alpha
			-\nabla_{\nabla_XY}\alpha\bigr)(Z).
	\]
	再次利用无挠性及第 1 问，
	\[
		(\nabla^2\alpha)(X,Y;Z)-(\nabla^2\alpha)(Y,X;Z)
		=-\alpha(R(X,Y)Z).
	\]
	在半径 \(a\) 的圆球面取一点处的正交单位向量 \(e_1,e_2\)，并取局部一形式使 \(\alpha_p=e_1^\flat\)。该点有
	\[
		R(e_1,e_2)e_2=a^{-2}e_1,
	\]
	所以当 \((X,Y;Z)=(e_1,e_2;e_2)\) 时，上述交换差为 \(-a^{-2}\ne0\)。公式的两侧都是张量，任意光滑局部延拓都给出相同的点处数值。
\end{solution}

\begin{exercise}\label{g3:ex:three-dimensional-ricci-positivity}
	在三维黎曼流形的一点 \(p\)，设 Ricci 张量对应的自伴算子具有特征值 \(r_1,r_2,r_3\)。
	\begin{enumerate}
		\item 利用 \cref{g3:prop:low-dimensional-curvature}，求曲率算子的三个特征值及一组相应的特征二向量。
		\item 证明该点所有截面曲率严格为正，当且仅当
		\[
			r_1+r_2>r_3,\qquad
			r_1+r_3>r_2,\qquad
			r_2+r_3>r_1.
		\]
		\item 给出三个严格正的 \(r_i\)，使某个截面曲率为负，并说明这些数据可以由某个局部黎曼度量实现。
	\end{enumerate}
\end{exercise}
\begin{solution}
	1. 取 \(g_p\) 正交单位基 \(e_1,e_2,e_3\)，使 \(\operatorname{Ric}(e_i,e_j)=r_i\delta_{ij}\)，并记 \(s=r_1+r_2+r_3\)。三维公式为
	\[
		\mathrm{Rm}=g\wedge\operatorname{Ric}-\frac{s}{4}g\wedge g.
	\]
	展开乘积，在此基中得到
	\[
		R_{ijkl}
		=\left(r_i+r_j-\frac{s}{2}\right)
			(\delta_{ik}\delta_{jl}-\delta_{il}\delta_{jk}).
	\]
	上式中的 \(i,j\) 是自由指标，不对它们求和。因此在正交单位二向量基
	\(e_1\wedge e_2,e_1\wedge e_3,e_2\wedge e_3\) 中，曲率算子为对角矩阵，三个特征值分别为
	\[
		\lambda_{12}=\frac{r_1+r_2-r_3}{2},\qquad
		\lambda_{13}=\frac{r_1+r_3-r_2}{2},\qquad
		\lambda_{23}=\frac{r_2+r_3-r_1}{2}.
	\]

	2. 若所有截面曲率为正，对由 \(e_i,e_j\) 张成的三个截面计算，便有 \(\lambda_{ij}>0\)，这给出所列三个不等式。反过来，若它们成立，则任意二维平面的单位二向量可写成
	\[
		\xi=b_{12}e_1\wedge e_2+b_{13}e_1\wedge e_3
			+b_{23}e_2\wedge e_3,\qquad
		b_{12}^2+b_{13}^2+b_{23}^2=1.
	\]
	于是该平面的截面曲率为
	\[
		\langle\mathcal R\xi,\xi\rangle
		=\lambda_{12}b_{12}^2+\lambda_{13}b_{13}^2
			+\lambda_{23}b_{23}^2>0.
	\]
	这证明充分性。三个不等式两两相加还给出每个 \(r_i>0\)，符合正截面曲率蕴含正 Ricci 曲率的收缩关系。

	3. 取 \((r_1,r_2,r_3)=(1,1,3)\)，则
	\[
		(\lambda_{12},\lambda_{13},\lambda_{23})
		=\left(-\frac12,\frac32,\frac32\right).
	\]
	Ricci 张量正定，但 \(e_1,e_2\) 张成的截面曲率为负。为了实现这些数值，在 \(\mathbb R^3\) 的 Euclid 内积 \(g_0\) 下取对称张量 \(h=\operatorname{diag}(1,1,3)\)，令
	\[
		A=g_0\wedge h-\frac54g_0\wedge g_0.
	\]
	两个乘积均有代数曲率对称性，故 \(A\) 是代数曲率张量。由收缩公式，
	\[
		\operatorname{Ric}(A)
		=\bigl(h+5g_0\bigr)-\frac54(4g_0)=h.
	\]
	因此 \cref{g3:prop:curvature-local-realization} 给出一个原点邻域上的光滑度量，其原点曲率正是 \(A\)。由连续性再缩小邻域，可使 Ricci 张量在整个邻域保持正定，同时坐标前两个方向张成的截面仍有负曲率。这里 Ricci 的正定性有严格的正下界，负截面值也有严格的负余量，所以二者都在小邻域内保留。
\end{solution}

\begin{exercise}\label{g3:ex:curvature-tensor-operator-norm}
	设 \((M^n,g)\) 为光滑黎曼流形，\(n\ge2\)。曲率张量范数按四个有序指标的全部分量平方和定义，曲率算子的平方范数记为 \(|\mathcal R|_{\mathrm{HS}}^2\)，定义为 \(\Lambda^2\mathrm{T}_pM\) 的一组正交单位基在该算子下的像向量长度平方和。
	\begin{enumerate}
		\item 证明
		\[
			|\mathrm{Rm}|^2=4|\mathcal R|_{\mathrm{HS}}^2.
		\]
		\item 结合式\eqref{g3:eq:curvature-norm-decomposition}与低维公式，证明
		\[
			n=3:\quad|\mathrm{Rm}|^2=4|\operatorname{Ric}|^2
				-\operatorname{Scal}^2,
			\qquad
			n=2:\quad|\mathrm{Rm}|^2=\operatorname{Scal}^2.
		\]
		解释第 1 问中系数 \(4\) 的来源。
	\end{enumerate}
\end{exercise}
\begin{solution}
	1. 取一点处的正交单位基 \(e_1,\ldots,e_n\)。由曲率算子的分量表达，
	\[
		|\mathcal R|_{\mathrm{HS}}^2
		=\sum_{i<j}|\mathcal R(e_i\wedge e_j)|^2
		=\sum_{i<j,\ k<l}R_{ijkl}^2.
	\]
	另一方面，
	\[
		|\mathrm{Rm}|^2=\sum_{i,j,k,l}R_{ijkl}^2.
	\]
	第一对或第二对指标相等时，分量为零。对每个 \(i<j,k<l\)，独立交换前一对或后一对，只改变分量的符号，故其平方在全部有序指标的和中出现四次。因此
	\[
		\sum_{i,j,k,l}R_{ijkl}^2
		=4\sum_{i<j,\ k<l}R_{ijkl}^2,
	\]
	得到所求等式。

	2. 三维时 \(W=0\)。记 \(s=\operatorname{Scal}\)、\(E=\operatorname{Ric}-\tfrac{s}{3}g\)。由于 \(\operatorname{tr}_gE=0\)、\(|g|^2=3\)，有
	\[
		|\operatorname{Ric}|^2=|E|^2+\frac{s^2}{3}.
	\]
	把 \(n=3\) 代入范数分解式，得到
	\[
		|\mathrm{Rm}|^2
		=4|E|^2+\frac{s^2}{3}
		=4|\operatorname{Ric}|^2-s^2.
	\]
	二维时曲率算子作用在一维空间上，等于 \(K\operatorname{id}\)，故
	\[
		|\mathrm{Rm}|^2=4K^2=(2K)^2=\operatorname{Scal}^2.
	\]
	也可直接列出四个可能非零的分量：
	\(R_{1212}=R_{2121}=K\)、\(R_{1221}=R_{2112}=-K\)。

	系数 \(4\) 来自两对反对称指标各自的两次计数。曲率算子使用的单位基是 \(e_i\wedge e_j\)、\(i<j\)，每个方向只列一次；张量范数则对所有四个有序指标求和。固定这一规范后，同一几何对象的张量表达与算子表达具有上述确定的范数换算。
\end{solution}

\begin{exercise}\label{g3:ex:flat-nonconstant-conformal-metric}
	设 \(U\subset\mathbb R^2\) 为开集，\(u\in C^\infty(U)\)，取度量
	\[
		g=e^{2u(x,y)}(\mathrm{d}x^2+\mathrm{d}y^2).
	\]
	\begin{enumerate}
		\item 用正交余标架的结构方程计算 Gauss 曲率，证明
		\[
			K=-e^{-2u}(u_{xx}+u_{yy}),
		\]
		并核对它与式\eqref{g3:eq:conformal-scalar}一致。
		\item 在整个 \(\mathbb R^2\) 取 \(u=ax\)，其中 \(a\ne0\) 为常数。证明度量平坦，并证明
		\[
			F(x,y)=\frac{e^{ax}}{a}\bigl(\cos(ay),\sin(ay)\bigr)
		\]
		给出每点附近的平直坐标。说明 \(F\) 为局部等距映射，但在整个 \(\mathbb R^2\) 上不是单射。
	\end{enumerate}
\end{exercise}
\begin{solution}
	1. 取正交单位余标架
	\[
		\theta^1=e^u\mathrm{d}x,\qquad
		\theta^2=e^u\mathrm{d}y.
	\]
	直接求外微分，
	\[
		\mathrm{d}\theta^1=-e^uu_y\,\mathrm{d}x\wedge\mathrm{d}y,
		\qquad
		\mathrm{d}\theta^2=e^uu_x\,\mathrm{d}x\wedge\mathrm{d}y.
	\]
	记 \(\alpha=\omega^2{}_1=-\omega^1{}_2\)。无挠结构方程为
	\[
		\mathrm{d}\theta^1-\alpha\wedge\theta^2=0,\qquad
		\mathrm{d}\theta^2+\alpha\wedge\theta^1=0.
	\]
	设 \(\alpha=A\,\mathrm{d}x+B\,\mathrm{d}y\)，第一式给出 \(A=-u_y\)，第二式给出 \(B=u_x\)。所以
	\[
		\alpha=-u_y\,\mathrm{d}x+u_x\,\mathrm{d}y,\qquad
		\mathrm{d}\alpha=(u_{xx}+u_{yy})\,\mathrm{d}x\wedge\mathrm{d}y.
	\]
	由式\eqref{g3:eq:surface-curvature-form}，
	\[
		\mathrm{d}\alpha=-K\theta^1\wedge\theta^2
		=-Ke^{2u}\,\mathrm{d}x\wedge\mathrm{d}y,
	\]
	比较系数即得 \(K=-e^{-2u}(u_{xx}+u_{yy})\)。

	为核对共形标量公式，底度量取 Euclid 度量 \(g_0\)。本书符号约定给出
	\[
		\Delta_{g_0}u=-(u_{xx}+u_{yy}),\qquad
		\operatorname{Scal}_{g_0}=0.
	\]
	二维公式因而为
	\[
		\operatorname{Scal}_g
		=2e^{-2u}\Delta_{g_0}u
		=-2e^{-2u}(u_{xx}+u_{yy})=2K,
	\]
	与结构方程计算相同。

	2. \(u=ax\) 的二阶偏导均为零，故 \(K=0\)。二维的全部曲率由 \(K\) 决定，所以此度量平坦，而 \(e^{2ax}\) 随 \(x\) 非恒定。

	记 \(F=(F^1,F^2)\)。对给定表达式求导得到
	\[
		\begin{aligned}
			\mathrm{d}F^1
			&=e^{ax}\bigl(\cos(ay)\,\mathrm{d}x-\sin(ay)\,\mathrm{d}y\bigr),\\
			\mathrm{d}F^2
			&=e^{ax}\bigl(\sin(ay)\,\mathrm{d}x+\cos(ay)\,\mathrm{d}y\bigr).
		\end{aligned}
	\]
	因此
	\[
		F^*(\mathrm{d}X^2+\mathrm{d}Y^2)
		=(\mathrm{d}F^1)^2+(\mathrm{d}F^2)^2
		=e^{2ax}(\mathrm{d}x^2+\mathrm{d}y^2)=g.
	\]
	导数矩阵的行列式为 \(e^{2ax}>0\)，所以每一点附近 \(F\) 都是到一个平面开集的微分同胚。将邻域限制到这个可逆范围后，\((X,Y)=(F^1,F^2)\) 正是度量系数为单位矩阵的平直坐标。

	但周期性给出
	\[
		F\left(x,y+\frac{2\pi}{|a|}\right)=F(x,y).
	\]
	右边的两个原像点不同，所以整个映射不单射。这一例子同时给出了非恒定共形因子的平坦度量及显式局部等距坐标，并表明这些坐标不能由同一映射直接拼成全局单射坐标。
\end{solution}

\begin{exercise}\label{g3:ex:scalar-flat-warped-model}
	固定 \(b>0\)，令
	\[
		r(t)=b(t+\sinh t\cosh t),\qquad f(r(t))=b\cosh^2t.
	\]
	证明这在整个实轴上定义了光滑正函数 \(f\)。在 \(M=\mathbb R\times S^2\) 上取 \(g=\mathrm dr^2+f(r)^2g_{S^2}\)，其中 \(S^2\) 的截面曲率为一。计算径向与纤维截面曲率、Ricci 特征值及标量曲率，并证明 \(g\) 完备。这个模型是否平坦？
\end{exercise}
\begin{solution}
	因为 \(r'(t)=2b\cosh^2t>0\)，且 \(r(t)\to\pm\infty\) 当 \(t\to\pm\infty\)，映射 \(r\) 是实轴的光滑微分同胚，所定义的 \(f\) 光滑且 \(f\ge b\)。按 \(r\) 求导得
	\[
		f'=\tanh t,\qquad
		f''=\frac{1}{2b\cosh^4t}=\frac{b}{2f^2},\qquad
		1-(f')^2=\frac bf.
	\]
	由式\eqref{g3:eq:warped-sectional-curvatures}，包含单位径向量 \(T=\partial_r\) 的截面与纯纤维截面分别满足
	\[
		K_{\mathrm{rad}}=-\frac{b}{2f^3},\qquad
		K_{\mathrm{tan}}=\frac{b}{f^3}.
	\]
	因此在 \(g\) 正交单位基 \((T,E_1,E_2)\) 中，Ricci 的特征值为
	\[
		-\frac{b}{f^3},\quad \frac{b}{2f^3},\quad \frac{b}{2f^3}.
	\]
	它们之和为零，故 \(\operatorname{Scal}_g=0\)，而径向截面处处具有负曲率，度量并不平坦。这里 \(f(1-(f')^2)=b\) 也可直接作为标量曲率为零的方程 \(2ff''=1-(f')^2\) 的第一积分核对；求导时不需要除以可能为零的 \(f'\)。

	最后证明完备。任意连接 \((r,x)\) 与 \((s,y)\) 的曲线满足
	\[
		L_g\ge |r-s|,\qquad L_g\ge b\,d_{S^2}(x,y).
	\]
	所以一个 \(d_g\)-Cauchy 序列的两个坐标分别在 \(\mathbb R\) 和 \(S^2\) 中是 Cauchy 序列，收敛到某个 \((r_\infty,x_\infty)\)。在一个包含所有充分靠后实坐标的紧区间上，\(f\) 有上界 \(C\)。先沿实方向、再沿球面方向连接，得到
	\[
		d_g((r,x),(r_\infty,x_\infty))
		\le |r-r_\infty|+C\,d_{S^2}(x,x_\infty).
	\]
	故原序列在 \(d_g\) 中收敛，证明了度量完备。这给出一个完备、标量曲率恒零而截面曲率同时有正有负的三维模型。
\end{solution}

\begin{exercise}\label{g3:ex:product-weyl-criterion}
	设 \((M^p,g_1)\)、\((N^q,g_2)\) 的截面曲率分别恒为 \(k\)、\(l\)，其中 \(p,q\ge2\)。在乘积上取 \(g=g_1\oplus g_2\)，记 \(n=p+q\)。
	\begin{enumerate}
		\item 求 \(\operatorname{Ric}_g=\lambda g\) 对 \(k,l\) 的充要条件。
		\item 证明 \(W_g=0\) 当且仅当 \(k+l=0\)。若同时满足第一问的条件，能得到什么结论？
	\end{enumerate}
\end{exercise}
\begin{solution}
	1. 由乘积曲率公式，两因子切方向上的 Ricci 分别为 \((p-1)k g_1\)、\((q-1)l g_2\)，混合分量为零。因此所求条件是
	\[
		(p-1)k=(q-1)l=\lambda.
	\]

	2. \(n\ge4\)，可使用 \cref{g3:thm:weyl-decomposition}。记标量曲率为
	\[
		s=p(p-1)k+q(q-1)l,
	\]
	则 \(A=(\operatorname{Ric}-s g/(2(n-1)))/(n-2)\) 在两个因子上分别是 \(a g_1\)、\(c g_2\)，其中
	\[
		a=\frac{(p-1)k-s/(2(n-1))}{n-2},\qquad
		c=\frac{(q-1)l-s/(2(n-1))}{n-2}.
	\]
	在第一因子的单位二维截面、第二因子的单位二维截面及混合单位截面上，\(g\wedge A\) 的值分别是 \(2a,2c,a+c\)。若 \(W=0\)，与乘积曲率比较便有
	\[
		k=2a,\qquad l=2c,\qquad a+c=0,
	\]
	从而 \(k+l=0\)。条件 \(p,q\ge2\) 保证前两个截面确实存在。

	反过来，若 \(l=-k\)，则
	\[
		s=(p-q)(n-1)k,\qquad a=\frac k2,\quad c=-\frac k2.
	\]
	在适合乘积分解的正交单位基中，\(\mathrm{Rm}\) 和 \(g\wedge A\) 的纯因子分量分别相等，所有混合分量也同时为零。因此四阶张量相等，\(W=0\)。若还要求第一问成立，则 \((p-1)k=-(q-1)k\)，即 \((n-2)k=0\)。故 \(k=l=0\)，乘积平坦。
\end{solution}

\begin{exercise}\label{g3:ex:conformal-scalar-sign}
	设 \((M^n,g)\) 紧、连通、无边界，\(n\ge3\)，且 \(\operatorname{Scal}_g=0\)。对正光滑函数 \(v\) 取 \(\widetilde g=v^{4/(n-2)}g\)。证明：若 \(\operatorname{Scal}_{\widetilde g}\) 处处非负或处处非正，则 \(v\) 为常数；因此任意非常数的这种共形变化都产生既取正值又取负值的标量曲率。
\end{exercise}
\begin{solution}
	由式\eqref{g3:eq:conformal-scalar-power}，
	\[
		\frac{4(n-1)}{n-2}\Delta_gv
		=\operatorname{Scal}_{\widetilde g}\,v^{(n+2)/(n-2)}.
	\]
	散度定理 \cref{g3:thm:divergence-theorem} 给出 \(\int_M\Delta_gv\,\mu_g=0\)。右端的权 \(v^{(n+2)/(n-2)}\) 严格为正；若标量曲率处处具有同一弱符号，其积分为零便迫使它处处为零。于是 \(\Delta_gv=0\)。再由第一章的能量公式，
	\[
		0=\int_Mv\Delta_gv\,\mu_g
		=\int_M|\operatorname{grad}_gv|_g^2\mu_g.
	\]
	所以 \(\mathrm dv=0\)，连通性使 \(v\) 为常数。若 \(v\) 非常数，标量曲率既不可能处处非负，也不可能处处非正，因而必有正值点和负值点。紧性和无边界条件在两次积分中都已使用。
\end{solution}

\begin{exercise}\label{g3:ex:conformal-halfspace-scalar-flat}
	在半空间 \(\{(x^1,\ldots,x^{n-1},y):y>0\}\subset\mathbb R^n\) 上，设 \(n\ge3\)，考虑
	\[
		\widetilde g=y^{4/(n-2)}\left(\sum_{\alpha=1}^{n-1}(\mathrm dx^\alpha)^2+\mathrm dy^2\right).
	\]
	计算其标量曲率、协变 Ricci 张量，以及径向和纯水平截面曲率。证明此度量共形平坦但不平坦，并判断它是否完备。
\end{exercise}
\begin{solution}
	记括号内的 Euclid 度量为 \(g_0\)，令 \(\alpha=2/(n-2)\)、\(u=\alpha\log y\)，则 \(\widetilde g=e^{2u}g_0\)。由于 \(v=y\) 满足 \(\Delta_{g_0}v=0\)，幂次形式的共形公式立即给出 \(\operatorname{Scal}_{\widetilde g}=0\)。又有
	\[
		\mathrm du=\frac\alpha y\mathrm dy,\qquad
		H_u=-\frac\alpha{y^2}\mathrm dy\otimes\mathrm dy,\qquad
		\Delta_{g_0}u=\frac\alpha{y^2}.
	\]
	代入 Ricci 变换公式并使用 \((n-2)\alpha=2\)，得到
	\[
		\widetilde{\operatorname{Ric}}
		=\frac{2}{(n-2)y^2}
		\bigl(n\,\mathrm dy\otimes\mathrm dy-g_0\bigr).
	\]
	这是坐标基中的协变张量，取迹时须使用 \(\widetilde g^{-1}=y^{-2\alpha}g_0^{-1}\)，不能直接把坐标系数当成曲率特征值。

	进一步，\(B_u=H_u-\mathrm du\otimes\mathrm du+\tfrac12|\mathrm du|^2g_0\) 在 Euclid 单位水平向量和 \(\partial_y\) 上的值分别为
	\[
		\frac{\alpha^2}{2y^2},\qquad
		-\frac{\alpha+\alpha^2/2}{y^2}.
	\]
	共形四阶公式给出：对一对 Euclid 正交单位向量，将它们再乘 \(e^{-u}\) 单位化后，新截面曲率是 \(-e^{-2u}\) 乘以这两个 \(B_u\) 对角值之和。因此
	\[
		K_{\mathrm{rad}}=\frac{2}{n-2}\,y^{-2n/(n-2)},\qquad
		K_{\mathrm{hor}}=-\frac{4}{(n-2)^2}\,y^{-2n/(n-2)}.
	\]
	这里径向截面包含 \(\partial_y\)，纯水平截面位于 \(y\) 的等值超平面；两类截面都存在。曲率非零且异号，而该度量按定义与平直度量共形，所以共形平坦。

	它并不完备。固定所有 \(x^\alpha=0\)，沿竖直方向向 \(y=0\) 行进的总长度为
	\[
		\int_0^1 y^{2/(n-2)}\,\mathrm dy=\frac{n-2}{n}<\infty.
	\]
	尤其点列 \((0,1/j)\) 的尾段距离被这段积分的尾部控制，因而是 Cauchy 序列。若它在流形中收敛，距离拓扑与流形拓扑一致便使其 \(y\) 坐标收敛到一个正数，与 \(1/j\to0\) 矛盾。故此例还说明，局部曲率和共形公式本身不保证完备性。
\end{solution}

\begin{exercise}\label{g3:ex:normalized-total-scalar-variation}
	设 \((M^n,g)\) 紧、连通、无边界，\(n\ge3\)。记
	\[
		V(g)=\int_M\mu_g,\qquad
		I(g)=\int_M\operatorname{Scal}_g\,\mu_g,\qquad
		\mathcal F(g)=V(g)^{-(n-2)/n}I(g).
	\]
	对 \(g_t=e^{2tu}g\)、\(u\in C^\infty(M)\)，计算 \(\left.\frac{\mathrm d}{\mathrm dt}\right|_{t=0}\mathcal F(g_t)\)。证明它对每个 \(u\) 都为零，当且仅当 \(g\) 的标量曲率为常数。
\end{exercise}
\begin{solution}
	记 \(s=\operatorname{Scal}_g\)。由共形标量曲率公式及密度的行列式公式，
	\[
		\left.\frac{\mathrm d}{\mathrm dt}\right|_0s_{g_t}
		=2(n-1)\Delta_gu-2us,\qquad
		\left.\frac{\mathrm d}{\mathrm dt}\right|_0\mu_{g_t}=nu\mu_g.
	\]
	紧性允许在积分号下求导，散度定理又使 Laplace 项的积分为零，故
	\[
		V'(0)=n\int_Mu\mu_g,\qquad
		I'(0)=(n-2)\int_Mus\mu_g.
	\]
	令 \(\bar s=I(g)/V(g)\)，乘积求导得到
	\[
		\mathcal F'(0)
		=(n-2)V(g)^{-(n-2)/n}
		\int_Mu(s-\bar s)\mu_g.
	\]
	若 \(s\) 为常数，则 \(s=\bar s\)，所有这些导数都为零。反过来，若对每个光滑 \(u\) 都为零，可取 \(u=s-\bar s\)。由于前面的系数为正，便有
	\[
		\int_M(s-\bar s)^2\mu_g=0,
	\]
	从而 \(s=\bar s\)。这一判据把常标量曲率条件表示成归一化总标量曲率沿全部共形方向的驻定条件。
\end{solution}

\begin{exercise}\label{g3:ex:ricci-invisible-local-curvature}
	在四维 Euclid 空间的正交单位基 \((e_1,e_2,e_3,e_4)\) 中，定义二向量空间上的对角自伴算子，其在基 \((e_i\wedge e_j)_{i<j}\) 上的特征值为
	\[
		k_{12}=1,\quad k_{13}=-1,\quad k_{14}=0,\quad
		k_{23}=0,\quad k_{24}=-1,\quad k_{34}=1.
	\]
	证明它对应一个非零的代数曲率张量 \(A\)，但 \(\operatorname{Ric}_A=0\)。将它实现为某个局部黎曼度量在原点的曲率，并说明该度量在原点的法坐标体积密度和小球体积为何在本章所求的首个曲率修正阶数上都与 Euclid 情形相同。
\end{exercise}
\begin{solution}
	令 \(A_{ijij}=k_{ij}\) 并按两对反对称性及成对交换对称性延拓，其余独立分量取零。这给出题中算子对应的四阶张量。第一 Bianchi 恒等式在四个指标互不相同时的各项都为零；有重复指标时，各项则由反对称性两两抵消。因此它确是代数曲率张量。

	Ricci 的非对角分量全为零，对角分量是每个方向相邻的三个截面值之和：
	\[
		\begin{aligned}
		(\operatorname{Ric}_A)_{11}&=1-1+0=0,&
		(\operatorname{Ric}_A)_{22}&=1+0-1=0,\\
		(\operatorname{Ric}_A)_{33}&=-1+0+1=0,&
		(\operatorname{Ric}_A)_{44}&=0-1+1=0.
		\end{aligned}
	\]
	但是 \(A_{1212}=1\)，所以 \(A\ne0\)。由 \cref{g3:prop:curvature-local-realization}，原点邻域中的正定度量
	\[
		g_{ij}(x)=\delta_{ij}-\frac13A_{ikjl}x^kx^l
	\]
	在原点的曲率恰是 \(A\)，因而那里 \(\operatorname{Ric}=0\)、\(\operatorname{Scal}=0\)，但 \(K(e_1,e_2)=1\)、\(K(e_1,e_3)=-1\)。

	对这个度量重新取原点的法坐标，由体积展开定理可得
	\[
		J_0(x)=1+O(|x|^3),\qquad
		\operatorname{Vol}_g(B_g(0,r))
		=\omega_4r^4\bigl(1+O(r^4)\bigr).
	\]
	因此密度的二次项及小球体积相对于 \(\omega_4r^4\) 的二次修正都消失。它们分别读取 Ricci 与标量曲率，不能据此推出整个曲率张量为零。本题只比较这两个指定阶数，没有断言更高阶项或精确体积也与平直情形相同。
\end{solution}
